% =====================================================================
% TÜRKÇE İNCELEME KOPYASI (gönderim sürümü değildir; SAUCIS yalnızca
% İngilizce kabul eder). Kaynak: ../manuscript/main.tex
% Olurluluğu Koruyan Güvenlik Katmanıyla V2G Şarjı için Boyuttan Bağımsız
% Güvenli Pekiştirmeli Öğrenme (çalışma başlığı, DD45; kararı yazar verir)
% Hedef dergi: SAUCIS (Sakarya University Journal of Computer and
% Information Sciences). Derleme: latexmk -pdf main_tr.tex
% Resmi şablon dosyası saucis.cls bu klasöre konursa otomatik olarak
% kullanılır; aksi hâlde standart article sınıfı kullanılır.
% =====================================================================
\newif\ifsaucis
\IfFileExists{saucis.cls}{\saucistrue\documentclass[final,9pt,onecolumn]{saucis}}{\saucisfalse\documentclass[10pt]{article}}

% ---------------- çift kör anahtar --------------------------------
% \anonymoustrue  -> hakem değerlendirmesi için anonim PDF (yazarlar,
%                   kurumlar, teşekkür adları ve depo bağlantısı gizlenir)
% \anonymousfalse -> yazar bilgilerini içeren PDF
\newif\ifanonymous
\anonymousfalse

% Türkçe inceleme kopyasının önsözü (manuscript_tr/main_tr.tex); İngilizce
% ana metin ve ek materyalle ortak önsözün uyarlanmış kopyasıdır.
% ---------------- compatibility preamble ----------------------------
\makeatletter
\ifsaucis\else
  \usepackage[a4paper,margin=2.5cm]{geometry}
  \usepackage[T1]{fontenc}
  \IfFileExists{lmodern.sty}{\usepackage{lmodern}}{}
  \setlength{\parindent}{0pt}          % SAUCIS: no paragraph indentation
  \setlength{\parskip}{0.5\baselineskip}
\fi
% Türkçe dil desteği (kısaltmalar kapalı: matematik ve tablolarda = ve : bozulmasın)
\usepackage[utf8]{inputenc}
\@ifpackageloaded{fontenc}{}{\usepackage[T1]{fontenc}}
\usepackage[turkish,shorthands=off]{babel}
\@ifpackageloaded{amsmath}{}{\usepackage{amsmath}}
\@ifpackageloaded{amssymb}{}{\usepackage{amssymb}}
\@ifpackageloaded{amsthm}{}{%
  \@ifundefined{proof}{}{\let\proof\relax\let\endproof\relax}%
  \usepackage{amsthm}}
\@ifpackageloaded{graphicx}{}{\usepackage{graphicx}}
\@ifpackageloaded{booktabs}{}{\usepackage{booktabs}}
\@ifpackageloaded{longtable}{}{\usepackage{longtable}}
\@ifpackageloaded{xcolor}{}{\usepackage{xcolor}}
\@ifpackageloaded{array}{}{\usepackage{array}}
\@ifpackageloaded{cite}{}{\usepackage{cite}}
\@ifpackageloaded{algorithm}{}{\usepackage{algorithm}}
\@ifpackageloaded{algpseudocode}{}{\@ifpackageloaded{algorithmic}{}{\usepackage{algpseudocode}}}
\@ifpackageloaded{caption}{}{\usepackage{caption}}
\IfFileExists{adjustbox.sty}{\@ifpackageloaded{adjustbox}{}{\usepackage{adjustbox}}}{}% generated result tables use adjustbox
\captionsetup{font=small,labelfont=bf}
\@ifundefined{floatname}{}{\floatname{algorithm}{Algoritma}}
\@ifpackageloaded{eurosym}{}{\IfFileExists{eurosym.sty}{\usepackage[gen]{eurosym}}{\providecommand{\euro}{\texteuro}}}
% cross-references between main text and supplementary material
\@ifpackageloaded{hyperref}{\usepackage{xr}}{\usepackage{xr-hyper}\usepackage[hidelinks]{hyperref}}
\makeatother

% Unicode symbols that generated tables may contain (pdflatex)
\makeatletter
\@ifundefined{DeclareUnicodeCharacter}{}{%
\DeclareUnicodeCharacter{03BB}{\ensuremath{\lambda}}\DeclareUnicodeCharacter{03A3}{\ensuremath{\Sigma}}%
\DeclareUnicodeCharacter{03C3}{\ensuremath{\sigma}}\DeclareUnicodeCharacter{03BC}{\ensuremath{\mu}}%
\DeclareUnicodeCharacter{03C1}{\ensuremath{\rho}}\DeclareUnicodeCharacter{03B1}{\ensuremath{\alpha}}%
\DeclareUnicodeCharacter{03B2}{\ensuremath{\beta}}\DeclareUnicodeCharacter{0394}{\ensuremath{\Delta}}%
\DeclareUnicodeCharacter{03B7}{\ensuremath{\eta}}\DeclareUnicodeCharacter{03C4}{\ensuremath{\tau}}%
\DeclareUnicodeCharacter{2264}{\ensuremath{\le}}\DeclareUnicodeCharacter{2265}{\ensuremath{\ge}}%
\DeclareUnicodeCharacter{2212}{\ensuremath{-}}\DeclareUnicodeCharacter{00B7}{\ensuremath{\cdot}}%
\DeclareUnicodeCharacter{00D7}{\ensuremath{\times}}\DeclareUnicodeCharacter{2248}{\ensuremath{\approx}}%
\DeclareUnicodeCharacter{2020}{\ensuremath{\dagger}}}
\makeatother
\setlength{\emergencystretch}{2em}
% anonymous builds: no creation date or trailer id in the PDF metadata
\ifanonymous\ifdefined\pdfinfoomitdate\pdfinfoomitdate=1\fi\ifdefined\pdftrailerid\pdftrailerid{}\fi\fi

% ---------------- theorem environments ------------------------------
\theoremstyle{plain}
\makeatletter
\@ifundefined{theorem}{\newtheorem{theorem}{Teorem}}{}
\@ifundefined{lemma}{\newtheorem{lemma}{Lemma}}{}
\@ifundefined{proposition}{\newtheorem{proposition}{Önerme}}{}
\@ifundefined{corollary}{\newtheorem{corollary}{Sonuç}}{}
\theoremstyle{definition}
\@ifundefined{assumption}{\newtheorem*{assumption}{Varsayım}}{}
\@ifundefined{remark}{\newtheorem{remark}{Not}
\@ifundefined{definition}{\newtheorem{definition}{Tanım}}{}}{}
\makeatother

% Türkçe başlık adları (babel'in varsayılanları: Özet, Kanıt)
\addto\captionsturkish{\renewcommand{\abstractname}{Öz}\renewcommand{\proofname}{İspat}%
  \renewcommand{\figurename}{Şekil}\renewcommand{\tablename}{Tablo}\renewcommand{\refname}{Kaynaklar}}

% ---------------- project macros ------------------------------------
\newcommand{\RES}[1]{\textcolor{red}{[SONUÇ: #1]}}
\newcommand{\TODO}[1]{\textcolor{blue}{[YAPILACAK: #1]}}
\newcommand{\method}{SA-Attn}          % method name; must match analysis/scripts/names.py
\newcommand{\Eq}[1]{Denklem~\ref{#1}}
\newcommand{\Eqs}[2]{Denklem~\ref{#1}--\ref{#2}}
\newcommand{\Fig}[1]{Şekil~\ref{#1}}
\newcommand{\Tab}[1]{Tablo~\ref{#1}}
\newcommand{\Sec}[1]{Bölüm~\ref{#1}}
\newcommand{\Alg}[1]{Algoritma~\ref{#1}}
\newcommand{\R}{\mathbb R}
\newcommand{\clip}{\operatorname{clip}}
\newcommand{\Ctil}{\tilde C}
\newcommand{\Crev}{C_{\mathrm{rev}}}
\newcommand{\Bh}{\hat B}
\newcommand{\Pnom}{\bar P^{\mathrm{nom}}}
\newcommand{\Pdis}{P^{\mathrm{dis}}}
\newcommand{\Ereq}{E^{\mathrm{req}}}
\newcommand{\pmn}{p^{\min}}
\newcommand{\sm}{\mathrm}

% figure from file, or a framed placeholder if the file is missing
\newcommand{\figorbox}[3]{%
  \IfFileExists{#1}{\includegraphics[width=#2]{#1}}{%
    \fbox{\parbox[c][#3][c]{0.9\linewidth}{\centering\small
      Şekil dosyası bulunamadı: \texttt{\detokenize{#1}}}}}}
% first existing file of three candidates (analysis output, analysis source, local copy)
\newcommand{\figchain}[4]{%
  \IfFileExists{#2}{\includegraphics[width=#1]{#2}}{%
  \IfFileExists{#3}{\includegraphics[width=#1]{#3}}{\figorbox{#4}{#1}{5cm}}}}
% Generated results. \analysisdir points to the analysis output tree of the
% planned paper/ + analysis/ layout; a build may override it before \input.
\providecommand{\analysisdir}{../../analysis}
% generated result table: analysis output, local copy in tables/, or placeholder
\newcommand{\resulttable}[2]{\IfFileExists{#1}{\input{#1}}{#2}}
\newcommand{\restable}[2]{%
  \IfFileExists{\analysisdir/tables/#1.tex}{\input{\analysisdir/tables/#1.tex}}{%
  \IfFileExists{tables/#1.tex}{\input{tables/#1.tex}}{#2}}}
% generated result figure: analysis output, local copy in figures/, or a red placeholder box
\newcommand{\resultfig}[2]{%
  \IfFileExists{\analysisdir/figures/#1}{\includegraphics[width=0.85\linewidth]{\analysisdir/figures/#1}}{%
  \IfFileExists{figures/#1}{\includegraphics[width=0.85\linewidth]{figures/#1}}{%
    \fbox{\parbox[c][5cm][c]{0.85\linewidth}{\centering\small
      \textcolor{red}{SONUÇ ŞEKLİ YER TUTUCUSU}\\[2mm]
      beklenen dosya: \texttt{\detokenize{#1}} (analysis/figures/)\\[1mm] #2}}}}}

\graphicspath{{figures/}}
\externaldocument[S-]{../supplementary_tr/supplementary_tr}

% =====================================================================
% Çalışma başlığı (DD45). Nihai başlığı yazar belirler.
\title{Olurluluğu Koruyan Güvenlik Katmanıyla V2G Şarjı için Boyuttan Bağımsız Güvenli Pekiştirmeli Öğrenme}
\ifanonymous
  \author{Anonim yazar(lar)\\ \small (yazar bilgileri çift kör değerlendirme için kaldırılmıştır)}
\else
  \author{Yazar Ad(lar)ı\thanks{Sorumlu yazar: \TODO{her yazarın adı, e-posta adresi ve ORCID kimliği}}\\
  \small \TODO{Kurum(lar), şehir, ülke}}
\fi
\date{\TODO{Nihai başlığı yazar belirler; bu, çalışma başlığıdır.}}

\begin{document}
\maketitle

\begin{center}
\fbox{\parbox{0.92\linewidth}{\small\itshape
Bu belge yazarın incelemesi için hazırlanmış Türkçe çeviridir; gönderim sürümü İngilizcedir. Şekillerin içindeki metinler ve analiz betiklerinin ürettiği sonuç tabloları (başlıkları dâhil) İngilizce bırakılmıştır; metindeki ondalık sayılar Türkçe yazımla (virgülle) verilmiş, matematik ifadelerindeki ve üretilen tablolardaki sayılar değiştirilmemiştir.}}
\end{center}

% ---------------------------------------------------------------------
% Öz (SAUCIS: <= 400 sözcük, atıf yok, şekil/tablo/bölüm göndermesi yok,
% üçüncü şahıs) ve anahtar kelimeler (3-5, alfabetik).
% Her sayı paper/results_number_sources.md dosyasında listelenmiştir.
% ---------------------------------------------------------------------
\begin{abstract}
Araçtan şebekeye (V2G) hizmetleri sunan şarj noktası işletmecileri, istasyonu trafo limiti içinde tutmak ve her sürücünün üzerinde anlaşılan enerjisini ayrılış anına kadar teslim etmek zorundadır.
Pekiştirmeli öğrenme (RL) tabanlı denetleyiciler bu iki gereksinimi genellikle ödüldeki cezalar olarak ele alır ve bu denetleyicilerin ağları tek bir istasyon boyutuna bağlıdır.
Bu makale, boyuttan bağımsız küme politikalarını bir güvenlik katmanıyla birleştirmektedir; bu katman, araç başına güç sınırları ile iki yönlü trafo limitinin oluşturduğu kümeye kesin bir izdüşümdür (projeksiyon).
Her elektrikli aracın (EV) alt sınırı, kabul edilen gereksinimini tam güçte erişilebilir tutar.
Herhangi bir ham eylem dizisi için trafo limitinin, simülatöre ve şebeke kapasite payına ilişkin koşullar altında her adımda sağlandığı ispatlanmıştır.
Kalış süresi boyunca olursuz adımla (trafonun tüm alt sınırları karşılayamadığı adımla) karşılaşmayan her EV, kabul edilen gereksinimini karşılar.

Dikkat tabanlı ve Deep Sets politikaları, EV2Gym'de 10 ve 20 şarj ünitesi üzerinde yakınsal politika optimizasyonu (PPO) ile eğitilmiş ve yeniden eğitim olmadan 200 şarj ünitesine kadar test edilmiştir.
Öğrenen her yöntem için iki eğitim tohumuyla, güvenlik katmanından geçirilerek çalıştırılan 102.750 test epizodunda hiç trafo aşırı yüklemesi oluşmamıştır; buna karşılık ceza ile eğitilen ve Lagrange tabanlı dikkat politikalarında sırasıyla 3.730 ve 3.274 aşırı yük adımı görülmüştür.
Ayrılışın bilindiği sınırı kullanan konformal olmayan her yöntemde, olursuz adım içermeyen hiçbir epizotta eksik enerjiyle ayrılan EV olmamıştır.

200 şarj ünitesinde dikkat tabanlı politikanın tam bilgili bir optimuma göre açığı şarj ünitesi ve epizot başına 0,90 avro olmuştur; 50 şarj ünitesinde eğitildiğinde bu değer 0,87 avrodur.
Deep Sets politikası 0,68 avroya ulaşmıştır; dolayısıyla dikkat mimarisi bir üstünlük göstermemiştir.
Aynı katmanla kullanılan bir fiyat eşiği kuralı, test ızgarasının her boyutunda dikkat tabanlı politikadan daha düşük ortalama açığa sahip olmuştur.
Deep Sets politikası bu kuralı yalnızca 20 şarj ünitesi ve üzerinde geçmiştir.

Öğrenen politikalar olursuz adımlara yol açmış; dikkat tabanlı politika için olursuz adımla karşılaşmayan EV'lerin payı 5 şarj ünitesinde \%65,9'dan 200 şarj ünitesinde \%92,6'ya yükselmiştir.
Ayrılış sınırıyla bu politika kabul edilen EV'lerin \%2,24'ünü gereksinimi karşılanmadan bırakmıştır; bu oran ceza ile eğitilen ikizinde \%1,86, sınır içermeyen bir varyantta ise \%0,99'dur.
Bu varyant her boyutta daha düşük bir açığa da sahip olmuştur.
Cezalı kâr açısından dikkat tabanlı politika, kabul edilen EV başına ikizinden 0,04 ile 3,46 avro daha yüksek puan almış; güven aralığı yalnızca 200 şarj ünitesinde sıfırı içermiştir.
\end{abstract}

\noindent\textbf{Anahtar Kelimeler:} Araçtan şebekeye, Elektrikli araç şarjı, Güvenli pekiştirmeli öğrenme, Güvenlik katmanı, Sıfır-atış aktarım

% ---------------------------------------------------------------------
\section{Giriş}\label{sec:intro}
% ---------------------------------------------------------------------

Bir şarj noktası işletmecisi (CPO), şarj üniteleri için enerjiyi zamanla değişen fiyatlardan satın alır; araçtan şebekeye (V2G) deşarj ise aynı işletmecinin, fiyatlar yüksekken park hâlindeki elektrikli araçlardan (EV) enerji satmasına olanak tanır.
İki katı gereksinim bu serbestliği sınırlar.
İstasyonun toplam gücü trafo limiti içinde kalmalı ve her EV, ayrılırken üzerinde anlaşılan enerjiye sahip olmalıdır.
Aksi hâlde işletmeci ya ekipmanı aşırı yükler ya da sürücüye verdiği hizmet taahhüdünü yerine getiremez.

Fiyatlar, varışlar ve ayrılışlar belirsiz olduğundan ve önceden modellenmeleri güç olduğundan pekiştirmeli öğrenme (RL), bu kontrol görevi için yaygın bir yaklaşım hâline gelmiştir \cite{qiu2023reinforcement,michailidis2025reinforcement}.
EV şarjına yönelik RL denetleyicilerinin çoğu, ayrılış enerji gereksinimini ödülde ağırlıklı bir ceza olarak kodlar \cite{sadeghianpourhamami2020definition,ye2022learning,korkas2024distributed,orfanoudakis2025ev2gym}.
Trafo limitini modelleyenler de aşırı yüklemeyi genellikle aynı şekilde cezalandırır \cite{orfanoudakis2025ev2gym,orfanoudakis2026graph}.
Ceza, ihlalleri kârla takas ettiğinden eğitilmiş bir politika test sırasında her iki gereksinimi de ihlal edebilir \cite{orfanoudakis2025ev2gym,tabas2022computationally}.
Kısıtlı formülasyonlar ise bunun yerine beklenen birikimli maliyeti sınırlar \cite{altman2021constrained}; ancak böyle bir sınır her adımda değil, epizotlar (episode) üzerinden ortalama olarak geçerlidir \cite{gu2024review}.

İkinci bir sınırlılık istasyon boyutuyla ilgilidir: EV2Gym simülatörünün varsayılan arayüzünde \cite{orfanoudakis2025ev2gym} olduğu gibi şarj ünitesi başına bir eylem içeren düz bir gözlem, eğitilmiş bir ağı tek bir şarj ünitesi sayısına bağlar.
Çizge ve dikkat tabanlı politikalar bu kısıtlamayı ortadan kaldırır \cite{orfanoudakis2025scalable,orfanoudakis2026graph,wang2026hetgps}; ne var ki boyuttan bağımsız bu denetleyicilerde kısıtlar hâlâ ödül terimleri, taklit verisi ya da sezgisel filtreler aracılığıyla devreye girer.

Güvenlik katmanları her eylemi uygulanmadan önce düzeltir ve limitleri yapıları gereği sağlayabilir \cite{dalal2018safe,chen2021enforcing,tabas2022computationally}.
EV istasyonları için mevcut katmanlar ya trafo kapasitesi ve şarj ünitesi akımı gibi adım başına limitleri kapsar \cite{wan2025logicinformed} ya da güvenli olmayan eylemleri başka eylemlerle değiştirir \cite{guan2024rule,zhang2025dumes}.
Yalnızca adım başına çalışan bir katman, ayrılış enerji gereksinimini korumaz; çünkü bu gereksinim, o anki eylemi kalış süresinin sonraki her adımıyla ilişkilendirir.
Araç başına en geç başlangıç sınırları bu bağlantıyı esneklik payına (laxity) dayalı çizelgelemede \cite{chen2022smoothed} ve yakın tarihli hiyerarşik bir V2G çizelgeleyicisinde \cite{yan2026hierarchical} ele almaktadır.
\Sec{sec:rw-position}, bu unsurların birleşimini en yakın çalışmalarla karşılaştırarak konumlandırmaktadır.

Bu makale, bir CPO istasyonunda V2G şarjı için boyuttan bağımsız güvenli kontrolü incelemektedir.
Bir küme politikası her şarj ünitesini bir belirteç olarak ele alır ve ortak ağırlıklarla belirteç başına bir ham güç eylemi üretir.
Ana gerçekleme olan güvenli dikkat tabanlı denetleyici (\method) öz-dikkat kullanır; bir Deep Sets politikası ise alternatif küme mimarisi işlevi görür.
Olurlu kümeye kesin bir izdüşüm (projeksiyon) olan güvenlik katmanı, ham eylemleri araç başına sınırlara ve trafo limitine uyan komut güçlerine dönüştürür.
Bir kutunun iki yönlü bir toplam kısıtıyla kesişimine izdüşüm klasik bir problemdir \cite{kiwiel2007breakpoint,wang2022efficient}; yumuşatılmış en az esneklik payı önce (sLLF) çizelgelemesinin hız kuralı da buna çok yakın bir biçime sahiptir \cite{chen2022smoothed}.
Olurluluk argümanımız, Chen ve arkadaşlarının esneklik payı argümanını şarj gücü azalmasına (taper), V2G deşarjına ve iki yönlü bir bağlantıya genişletmektedir.

Makalenin üç katkısı vardır.
\begin{enumerate}
\item Araç başına alt sınırı şarj gücü azalmasını içeren bir güvenlik katmanı türetilmekte ve Teorem~\ref{thm:invariance} bu katmanın sonucunu iki katmanda ifade etmektedir. Simülatöre ve şebeke kapasite payına ilişkin koşullar altında trafo limiti, her ham eylem dizisi için her adımda sağlanır; kabul edilen ayrılış gereksinimi ise kalış süresi boyunca olursuz adımla karşılaşmayan her EV tarafından karşılanır.
\item Boyuttan bağımsız iki küme politikası, bir dikkat tabanlı politika ve bir Deep Sets politikası, katman eğitim döngüsünün içindeyken 10 ve 20 şarj üniteli istasyonlarda eğitilmektedir. Her ikisi de yeniden eğitim olmadan 200 şarj ünitesine kadar, 50 şarj ünitesinde eğitilen bir referans politikaya karşı test edilmektedir.
\item EV2Gym'de trafo limitinin katmanlı ve katmansız durumda ve koşullu ayrılış sonucunun epizot düzeyinde ölçüme dayalı bir değerlendirmesi raporlanmaktadır. Değerlendirme ayrıca öğrenen politikaların ne sıklıkla olursuz adımlara yol açtığını, garantinin maliyetini, aynı katmanla kullanılan basit kuralları ve bölünmüş konformal bir ayrılış sınırını ölçmektedir.
\end{enumerate}
Sonuçlar, dikkat tabanlı politikanın Deep Sets politikasına üstünlüğünü göstermemektedir.
Güvenlik katmanıyla kullanılan bir fiyat eşiği kuralı, test ızgarasının her istasyon boyutunda dikkat tabanlı politikadan daha düşük ortalama açığa sahip olmuştur.
Bu kuralı yalnızca Deep Sets politikası, o da yalnızca 20 şarj ünitesi ve üzerinde geçmiştir.
Dikkat tabanlı politika için ayrılış sınırı, sınır içermeyen bir varyanta göre gereksinimi karşılanmayan daha fazla EV ile birlikte gelmiştir.
Bileşenlerin hiçbiri tek başına yeni değildir; dolayısıyla makalenin iddiası bunların birleşimine, iki katmanlı sonuca ve bu sonucun ölçülen davranışına dayanmaktadır.

\Sec{sec:related} ilgili çalışmaları gözden geçirmekte, \Sec{sec:problem} modeli kısmen gözlemlenebilir bir Markov karar süreci (POMDP) olarak ifade etmekte ve \Sec{sec:method} yöntemi sunmaktadır.
Bölüm~\ref{sec:setup}--\ref{sec:conclusion} sırasıyla deney düzeneğini, bulguları, tartışmayı ve sonucu kapsamakta; tüm ispatlar ve ek ayrıntılar ise ek materyalde yer almaktadır.

% ---------------------------------------------------------------------
\section{İlgili Çalışmalar}\label{sec:related}
% ---------------------------------------------------------------------

\subsection{EV şarjı ve V2G için pekiştirmeli öğrenme}\label{sec:rw-rl}

EV şarjına yönelik derin RL çalışmaları tek araçlı ajanlardan \cite{li2020constrained} fotovoltaik, depolama ve V2G içeren istasyonlara \cite{korkas2024distributed,zhang2025dumes}, koordineli istasyonlara \cite{song2025longterm} ve dağıtım şebekesi limitleri altındaki filolara \cite{zhang2024hamlet,caballe2026largescale} doğru ilerlemiştir.
İki derleme çalışması; güvenliği, simülasyona bağımlılığı ve standart bir kıyaslama ortamının eksikliğini açık problemler arasında saymaktadır \cite{qiu2023reinforcement,michailidis2025reinforcement}.
EV2Gym, şarj üniteleri esnek olmayan yük ve güneş üretimi bulunan trafoların arkasında yer alan bir CPO'yu modellemektedir \cite{orfanoudakis2025ev2gym}.
EV2Gym'in V2G ödülü trafo aşırı yüklemesini ve düşük kullanıcı memnuniyetini cezalandırır; V2G kıyaslamasında cezaya dayalı ajanlar kâr elde etmiş, ancak birçok kullanıcıyı talep ettikleri enerjinin altında bırakmıştır.
Adaptive Charging Network (ACN), katı altyapı limitleri ve en iyi çaba esaslı enerji teslimiyle kayan ufukta bir dışbükey program çözmektedir \cite{lee2021adaptive}.
EV2Gym içindeki ekonomik bir model öngörülü kontrol (MPC) yöntemi, ayrılış ve trafo kısıtlarını katı olarak uygulamaktadır \cite{diazlondono2026open}.

\subsection{Kısıt yönetimi ve güvenli pekiştirmeli öğrenme}\label{sec:rw-safe}

Ödüldeki ağırlıklı ihlal terimleri test sırasında ihlalleri önlemez \cite{tabas2022computationally,orfanoudakis2025ev2gym}.
Kısıtlı Markov karar süreci (CMDP) bunun yerine beklenen birikimli bir maliyeti sınırlar \cite{altman2021constrained} ve Lagrange \cite{tessler2018reward} ya da güven bölgesi \cite{achiam2017constrained} güncellemeleriyle çözülür.
Derleme çalışmaları, bu tür sınırların her adımda değil beklenen birikimli bir maliyet için geçerli olduğunu belirtmekte \cite{gu2024review} ve Lagrange yöntemlerini yumuşak garantiler olarak sınıflandırmaktadır \cite{su2025review}.
CMDP formülasyonları gerilim limitleri \cite{fan2024safety}, trafo ısıl zorlanması \cite{zhang2025modelfree,zhang2026safe} ve V2G toplayıcıları \cite{eland2026safe} ile EV şarjına uygulanmıştır.
Lagrange tabanlı çok ajanlı bir yöntem ayrıca her istasyonun yakın geçmişini bir Transformer ile kodlamaktadır \cite{huang2026safe}.

İkinci bir yöntem ailesi her eylemi uygulanmadan önce kalkanlar (shield) \cite{alshiekh2018safe}, öğrenilmiş doğrusal güvenlik katmanları \cite{dalal2018safe} ya da yedek denetleyicili kuadratik programlar \cite{ceusters2023adaptive} aracılığıyla düzeltir.
EV alanındaki örnekler arasında kural tabanlı kalkanlar \cite{guan2024rule}, eylem değiştirme katmanları \cite{zhang2025dumes}, şebeke yönetmeliğine dayalı kalkanlama \cite{demirel2025multiobjective} ve güç akışı kalkanları \cite{zaman2026graph} bulunmaktadır.
Böyle bir filtre ancak iç modeli kadar güvenlidir; HetGPS ön baskısının yazarları da doğrusallaştırılmış bir güç akışı izdüşümüyle birlikte kullanılan öğrenilmiş risk modeli için noktasal bir olurluluk iddiasında bulunmamaktadır \cite{wang2026hetgps}.
EV2Gym ekibinin üyeleri ayrıca çizge ile güçlendirilmiş bir ikiz gecikmeli derin deterministik politika gradyanı (TD3) ajanını \cite{fujimoto2018addressing} bir güvenlik katmanıyla birleştirmiştir \cite{orfanoudakis2026can}.
Bölüm özetine göre bu katmanın amacı şarj gereksinimlerini korumak ve güç limitlerini uygulamaktır.

Üçüncü bir aile, politikanın çıktısını yapısı gereği olurlu kılar.
Türevlenebilir dışbükey optimizasyon katmanları gradyanları kesin bir çözümden geçirir \cite{amos2017optnet,agrawal2019differentiable}; Chen ve arkadaşları da böyle bir izdüşümü enerji sistemi politikalarının içinde kullanmıştır \cite{chen2021enforcing}.
Gauge eşlemeleri doğrudan olurlu noktalar verir \cite{li2023learning}; Tabas ve Zhang, öğrenilmiş bir frekans denetleyicisinin kontrollü değişmez bir güvenli kümede kaldığını ispatlamak için bu eşlemelerden birini kullanmıştır \cite{tabas2022computationally}.
EV istasyonları için Wan ve arkadaşlarının bir ön baskısı, trafo ve şarj ünitesi akımı limitlerini kapsayan ancak ayrılış enerjisini kapsamayan bir kuadratik program izdüşümünün adım başına olurluluğunu ispatlamaktadır \cite{wan2025logicinformed}.
Yan ve arkadaşlarının bir ön baskısı araç başına sınırları kalan enerji ve süreden türetmektedir; bu çalışmanın biçimsel sonucu birinci dereceden güç dengesi olurluluğuyla ilgilidir \cite{yan2026hierarchical}.
İzdüşüm öğrenmeyi de etkiler; çünkü naif izdüşüm politika gradyanlarını yanlı hâle getirir \cite{gros2020safe} ve aktif bir izdüşümün gradyanları sıfırdır \cite{kasaura2023benchmarking}.
Stokastik politikalar için yakın tarihli bir ön baskı, ortamın içindeki izdüşümün bir politika katmanı olarak izdüşüme denk olduğunu göstermektedir \cite{markgraf2025safe}.

\subsection{Boyuttan bağımsız mimariler}\label{sec:rw-scale}

Deep Sets \cite{zaheer2017deep}, öz-dikkat \cite{vaswani2017attention} ve Set Transformer \cite{lee2019set}, ortak ağırlıklarla her boyuttaki kümeleri işleyebilir; UPDeT bu fikri farklı boyutlardaki çok ajanlı görevlere uygulamıştır \cite{hu2021updet}.
Boyut aktarımı kendiliğinden gerçekleşmez; küçük çizgeler üzerinde eğitilen sinirsel rotalama çözücüleri daha büyük çizgelerde sıklıkla başarısız olur \cite{joshi2022learning}.
EV şarjında boyuttan bağımsız ilk durum gösterimleri toplulaştırmadan elde edilmiş; Sadeghianpourhamami ve arkadaşları ayrıca ölçeklenmiş varış hızları altında mükemmel bilgiye sahip bir optimuma olan açığı raporlamıştır \cite{sadeghianpourhamami2020definition}.
Ancak toplulaştırılmış gösterimler araç başına ayrıntıyı gizler ya da bu ayrıntıyı ayrıştırma kurallarına taşır \cite{zhang2022transfer,salesortiz2024efficient}.

Heterojen çizge politikasının (EV-GNN) ön baskı sürümü, her ölçek için ayrı bir model kullanarak ve yapısı gereği hiçbir kısıtı uygulamadan 1.000'e kadar şarj noktasını koordine etmektedir \cite{orfanoudakis2025scalable}.
Buna karşılık, 25 şarj ünitesinde eğitilen çizge ile güçlendirilmiş karar Transformer'ı (GNN-DT), yeniden eğitim olmadan 5 ile 100 arasında şarj ünitesinde test edilmiş ve boyut büyüdükçe performansı düşmüştür \cite{orfanoudakis2026graph}.
1.236 araçla eğitilen bir HetGPS politikası, yeniden eğitim olmadan 2.239 ve 3.218 araçlı fiderlere aktarılabilmektedir \cite{wang2026hetgps}.
Benzer şekilde, gerilim kontrolüne yönelik çizge tabanlı bir aktör, simülasyonda eğitimdekinden daha büyük bir fider üzerinde değerlendirilmiştir \cite{jin2026telemetry}.

Boyuttan bağımsız bu denetleyicilerde kısıtlar ödül terimleri, taklit verisi ya da sezgisel filtreler aracılığıyla devreye girer.

\subsection{Esneklik payına dayalı çizelgeleme ve ayrılış belirsizliği}\label{sec:rw-laxity}

Chen ve arkadaşları, bir çizelgenin ancak ve ancak her esneklik payı, yani kalan süre ile şarjı tamamlamak için hâlâ gereken süre arasındaki fark, negatif olmadığında olurlu olduğunu ispatlamıştır \cite{chen2022smoothed}.
Onların sLLF kuralı, istasyon güç limiti altında her hızı ortak bir düzey etrafında kırpar; ortaya koydukları sınırlar ise çevrim içi olurluluğun genellikle ek güç gerektirdiğine işaret etmektedir.
Bir ön baskı, katı son tarihlerin ve bir trafo limitinin tıkanıklık altında birlikte olursuz hâle gelebileceğini göstermektedir \cite{reddy2026deadline}; en az esneklik payı önce (LLF) öncelikleri ise kablo ve trafo limitleri altında iyi performans göstermektedir \cite{vanhuffelen2025grid,christensen2025multiagent}.
Bu yöntemler, kullanıcıların ayrılış zamanını ve enerjiyi varışta bildirdiğini varsayar.
Bağlı kalma süreleri zor tahmin edilir \cite{dyo2022behavior}, kullanıcılar son tarihlerini olduğundan erken bildirerek kazanç sağlayabilir \cite{ferragut2022scheduling} ve ayrılış tahminleri önceliklendirmeyi iyileştirir \cite{frendo2021improving}.
Yakın tarihli bir ön baskı, bir ayrılış tahmincisini şarj politikasıyla uçtan uca eğitmekte, ancak karşılanmayan enerjiyi yumuşak bir ceza olarak tutmakta ve herhangi bir kapsama ifadesi sunmamaktadır \cite{gabriele2026forecasting}.
Konformal tahmin, değiştirilebilirlik (exchangeability) altında dağılımdan bağımsız kapsama sağlar \cite{lei2018distributionfree,angelopoulos2023conformal} ve güvenli planlamada kısıtları sıkılaştırmak için kullanılmıştır \cite{lindemann2023safe}.

\subsection{Konumlandırma}\label{sec:rw-position}

\Tab{tab:positioning} incelendiğinde görüldüğü gibi, tasarımımızın her bileşeni bir biçimde zaten mevcuttur.
EV şarjı için sıfır-atış aktarıma sahip boyuttan bağımsız politikalar \cite{orfanoudakis2026graph,wang2026hetgps}, güvenli RL'de ispatlanabilir eylem düzeltme katmanları \cite{tabas2022computationally,wan2025logicinformed} bulunmaktadır; ortak bir düzey etrafında kırpmaya dayanan çok yakın bir kural da sLLF'de yer almaktadır \cite{chen2022smoothed}.
Araç başına en geç başlangıç sınırları Yan ve arkadaşlarının çizelgeleyicisinde \cite{yan2026hierarchical} görülmekte; küçükte eğitip büyükte test etmeye ve bir optimuma olan açığa dayalı değerlendirme ise en azından \cite{sadeghianpourhamami2020definition} çalışmasına kadar uzanmaktadır.
Yakın tarihli ve konuyla çok yakından ilişkili üç çalışmaya yalnızca özetleri üzerinden erişilebilmiştir \cite{orfanoudakis2026can,song2025longterm,eland2026safe}; bu çalışmalar yalnızca bu düzeyde aktarılmaktadır.
Bildiğimiz kadarıyla, boyuttan bağımsız öğrenilmiş bir politikayı, adım başına bir trafo garantisini ve çok adımlı bir ayrılış argümanını tek bir denetleyicide birleştiren önceki bir çalışma bulunmamaktadır.
% Yazar notu (basılmaz): gönderimden önce orfanoudakis2026can, song2025longterm, eland2026safe ve yayımlanmış EV-GNN makalesini tam metin üzerinden yeniden kontrol edin (boşluk analizi R1-R4).

\begin{table}[t]
\centering
\caption{En yakın çalışmalara göre konumlandırma. ``Ceza'' bir ödül terimini; ``katı'' ise yapısı gereği uygulanan bir kısıtı ifade eder.}
\label{tab:positioning}
\footnotesize
\setlength{\tabcolsep}{3pt}
\begin{tabular}{@{}>{\raggedright}p{2.6cm}>{\raggedright}p{2.0cm}>{\raggedright}p{2.0cm}>{\raggedright}p{2.0cm}>{\raggedright}p{2.2cm}>{\raggedright}p{1.9cm}>{\raggedright\arraybackslash}p{2.0cm}@{}}
\toprule
Çalışma & Boyuttan bağımsız politika & Ayrılış enerjisi & Trafo limiti & Garanti & Sıfır-atış boyut testi & Optimuma olan açık \\
\midrule
EV2Gym ajanları \cite{orfanoudakis2025ev2gym} & hayır (düz) & ceza & ceza & yok & hayır & karşılaştırma yöntemi olarak optimum \\
EV-GNN, ön baskı sürümü \cite{orfanoudakis2025scalable} & evet (çizge) & yapısı gereği değil & yapısı gereği değil & yok & her ölçek için ayrı model & her ölçek için optimizasyon tabanlı karşılaştırma \\
GNN-DT \cite{orfanoudakis2026graph} & evet (çizge) & ceza & ceza & yok & 25'ten 5--100 şarj ünitesine & kâhin (oracle) düzeyine yakın ödül (dergi özeti) \\
HetGPS, ön baskı \cite{wang2026hetgps} & evet (çizge dikkati) & ödül ve sezgisel yöntem & kısıtlanmamış & deneysel (gerilim filtresi) & 1.236'dan 3.218 araca & hayır \\
Yan vd., ön baskı \cite{yan2026hierarchical} & hayır & katı (araç başına kural) & modellenmemiş & birinci dereceden güç dengesi & hayır & MPC'ye göre açık \\
Wan vd., ön baskı \cite{wan2025logicinformed} & hayır & modellenmemiş & katı (kuadratik program izdüşümü) & adım başına olurluluk & hayır & hayır \\
sLLF \cite{chen2022smoothed} & kural, her boyut & esneklik payı kuralı & katı (ortak düzey) & olurluluk teoremleri & öğrenilmemiş & hayır \\
Sadeghian\-pourhamami vd. \cite{sadeghianpourhamami2020definition} & evet (toplu durum) & ceza & modellenmemiş & yok & ölçeklenmiş varışlar & evet \\
Orfanoudakis vd. \cite{orfanoudakis2026can} (özet) & çizge politikası & güvenlik katmanı (belirtilen amaç) & güvenlik katmanı (belirtilen amaç) & özette belirtilmemiş & özette belirtilmemiş & özette belirtilmemiş \\
\midrule
Bu çalışma & evet (şarj ünitesi belirteçleri üzerinde dikkat ve Deep Sets) & araç başına sınır, koşullu (Teo.~1b) & katı (ortak kaydırma) & trafo: her adımda (A3, A6); ayrılış: araç başına, koşullu & 10/20'den 200 şarj ünitesine & yumuşak ve katı kısıtlı optimumlar \\
\bottomrule
\end{tabular}
\end{table}

% ---------------------------------------------------------------------
\section{Problem Formülasyonu}\label{sec:problem}
% ---------------------------------------------------------------------

\subsection{İstasyon, EV'ler ve enerji dinamiği}\label{sec:station}

Bir CPO, tek bir trafonun arkasında $N$ şarj ünitesine $\mathcal N=\{1,\dots,N\}$ sahip bir istasyon işletir; her şarj ünitesine en fazla bir EV bağlanır ve bu EV'ye ait büyüklükler şarj ünitesinin $i$ indisini taşır.
Zaman, $\Delta t$ [h] uzunluğundaki $t=0,\dots,T-1$ adımlarına bölünmüştür.
İstasyon her adımda trafo anma gücünü $P_{\sm{tr}}(t)$, esnek olmayan yükü $P^{\sm{inf}}(t)\ge0$ ve fotovoltaik (PV) üretimi $P^{\sm{pv}}(t)\ge0$ kW cinsinden gözlemler.
Enerji $c^{\sm{ch}}(t)$ fiyatından satın alınır ve $c^{\sm{dis}}(t)$ [\texteuro/kWh] fiyatından satılır; her iki fiyat da negatif olabilir.

EV $i$, $t_i^{\sm{arr}}$ adımında varır ve son kontrol edilen adım olan ayrılış adımı $d_i$'ye kadar kontrol edilir; dolayısıyla ayrılış enerjisi $E_i(d_i+1)$'dir; burada $E_i(t)$, $t$ adımının başındaki batarya enerjisidir.
$t$ adımında EV'nin $\tau_i(t)=d_i-t+1$ kalan eylemi vardır ve $\mathcal N(t)$, $n(t)$ adet dolu şarj ünitesinin kümesidir.
$p_i(t)$ [kW], $i$ şarj ünitesine komut edilen güç olsun; bu güç şarj için pozitif, V2G deşarjı için negatiftir. $\tilde p_i(t)$ ise simülatörün raporladığı güç olsun.
$p^{+}=\max(p,0)$ ve $p^{-}=\max(-p,0)$ olmak üzere batarya enerjisi
\begin{equation}\label{eq:dynamics}
E_i(t+1)=E_i(t)+\Delta t\,g_i\big(p_i(t)\big),\qquad g_i(p)=\eta_i^{\sm c}p^{+}-p^{-}/\eta_i^{\sm d},
\end{equation}
biçiminde değişir; burada $\eta_i^{\sm c},\eta_i^{\sm d}\in(0,1]$ şarj ve deşarj verimleridir.
Adım başına fiziksel limitler gücü
\begin{equation}\label{eq:limits}
\pmn_i(E)\le p_i(t)\le \bar P_i(E),\qquad \pmn_i(E)=-\min\Big\{\Pdis_i,\ \eta_i^{\sm d}\frac{(E-E_i^{\min})^{+}}{\Delta t}\Big\},\qquad E=E_i(t),
\end{equation}
aralığıyla sınırlar; burada en yüksek şarj gücü $\bar P_i(E)\ge0$ şarj gücü azalmasını ve kapasite limitini içerir, $\Pnom_i=\sup_E\bar P_i(E)$'dir ve $\Pdis_i$ en yüksek deşarj gücüdür.
Kuram genel $g_i$ eşlemesi için geçerlidir; ancak EV2Gym, batarya tarafındaki enerji değişimini şarj ünitesi gücü olarak raporlar, faturalandırır ve ölçer \cite{orfanoudakis2025ev2gym}.
Bu nedenle deneylerimizde $\eta_i^{\sm c}=\eta_i^{\sm d}=1$ alınmaktadır; bu kabul, V2G arbitrajındaki gidiş-dönüş kaybını ortadan kaldırmaktadır (\Sec{sec:limitations}).

\subsection{Ayrılış gereksinimi ve trafo limiti}\label{sec:requirement}

Her EV, istenen bir ayrılış enerjisi $E_i^{\sm{des}}$ talep eder ve istasyon düzeyindeki bir talep işleme hattı bu talebi kabul edilen gereksinim $\Ereq_i$'ye dönüştürür.
İşleme hattı, her test tekrarı (replay) için herhangi bir denetleyiciden önce bir kez çalışır; karşılaştırılan her yöntem ve her iki optimum aynı kabul edilen gereksinimi alır.
Varışta talep, kalış süresi boyunca tam güçte şarjın ulaşabileceği enerjiyle, bu değerin yaklaşık $10^{-6}$ kWh altında olacak şekilde sınırlandırılır.
Ardından kabul kontrolü küçük bir olurluluk doğrusal programı (LP) çözer ve bağlı EV'lerin tümüne öngörülen trafo limitleri altında hizmet verilemeyecekse yeni talebi düşürür.
Daha önce kabul edilmiş talepler hiçbir zaman azaltılmaz.
Dolayısıyla
\begin{equation}\label{eq:cap}
\Ereq_i\le\min\big\{E_i^{\sm{des}},\ \text{kalış süresince }E_i^{\sm{arr}}\text{ değerinden tam güçle erişilebilen enerji}\big\}
\end{equation}
olur; ayrıntılar Bölüm~\ref{S-sec:pipeline} içinde verilmektedir.
Ayrılış gereksinimi şudur:
\begin{equation}\label{eq:requirement}
E_i(d_i+1)\ \ge\ \Ereq_i .
\end{equation}
Bu sınırlama olmasaydı bazı varışlara hiçbir denetleyici hizmet veremezdi ve bunlara ilişkin bir garanti anlamsız olurdu.
Her yöntem için eksik enerji, ortak kabul edilen gereksinime göre ölçülür; işleme hattının çıkardığı enerji ise ayrıca raporlanır.

Net trafo gücü $P^{\sm{inf}}(t)-P^{\sm{pv}}(t)+\sum_{i}p_i(t)$, $[-P_{\sm{tr}}(t),P_{\sm{tr}}(t)]$ aralığında kalmalıdır; dolayısıyla EV toplamı için trafo limiti
\begin{equation}\label{eq:transformer}
-\Crev(t)\le\sum_{i\in\mathcal N}p_i(t)\le C(t),\qquad
C(t)=P_{\sm{tr}}(t)-P^{\sm{inf}}(t)+P^{\sm{pv}}(t),\qquad \Crev(t)=P_{\sm{tr}}(t)+P^{\sm{inf}}(t)-P^{\sm{pv}}(t)
\end{equation}
biçimindedir.
İleri yön limiti $C(t)$ şarjı, ters yön limiti $\Crev(t)$ ise V2G ile dışarı verilen gücü sınırlar.
Güvenlik katmanı mevcut adım için ölçülen yük ve PV değerlerini kullanırken politika gelecekteki değerlerin yalnızca gürültülü tahminlerini gözlemler.

\subsection{Çevrim içi kontrol problemi ve optimumlar}\label{sec:online}

CPO her adımda $p(t)\in\R^N$ değerini yalnızca $t$ adımında mevcut olan bilgiyi kullanarak seçer.
Bir $\pi$ denetleyicisinin bir epizottaki \texteuro{} cinsinden kârı
\begin{equation}\label{eq:objective}
J^{\pi}=\sum_{t=0}^{T-1}\Delta t\sum_{i\in\mathcal N(t)}\big(c^{\sm{dis}}(t)\,\tilde p_i^{-}(t)-c^{\sm{ch}}(t)\,\tilde p_i^{+}(t)\big)
\end{equation}
olarak tanımlanır ve denetleyici bunun beklenen değerini \Eqs{eq:dynamics}{eq:limits}, her EV için \Eq{eq:requirement} ve her adımda \Eq{eq:transformer} kısıtları altında en büyüklemelidir.
Her yöntem için $w_E=w_P=10$\,\texteuro/kWh ağırlıklarıyla cezalı kâr
\begin{equation}\label{eq:jpen}
J^{\sm{pen}}(\pi)=J^{\pi}-w_E\sum_i\Delta E_i^{\sm{sh}}-w_P\,\Delta t\sum_t\Big[\Big(\sum_i\tilde p_i(t)-C(t)\Big)^{+}+\Big(-\Crev(t)-\sum_i\tilde p_i(t)\Big)^{+}\Big]
\end{equation}
olarak tanımlanır; burada $\Delta E_i^{\sm{sh}}=(\Ereq_i-E_i(d_i+1))^{+}$ ayrılıştaki eksik enerjidir.

HiGHS \cite{huangfu2018parallelizing} ile çözülen iki tam bilgili LP, epizodun tamamını bilir ve referans olarak kullanılır.
Her ikisi de EV2Gym'in şarj gücü azalmasını dış doğrusal kısıt satırlarıyla modeller; bu satırlar tek adımlık bir geçiş bandı ve \%1'lik bir pay dışında kesindir (Lemma~\ref{S-lem:taperrows}); ayrıca her ikisi de EV2Gym'in faturalandırılmayan 0,01 kWh'lik yuvarlama kazancını hesaba katar.
Yumuşak cezalı optimum $J^{\sm{pen,LP}}$, eksik enerji ve aşırı yük için gevşek değişkenlerle \Eq{eq:jpen} ifadesini en büyükler.
Eksik enerji ortak kabul edilen gereksinime göre ölçüldüğünde, eksik enerji ya da aşırı yük içeren yörüngeler de dâhil olmak üzere simüle edilen her yörünge için $J^{\sm{pen}}(\pi)\le J^{\sm{pen,LP}}$ sağlanır (Önerme~\ref{S-prop:softoracle}).
Dolayısıyla ana ölçüt, şarj ünitesi ve epizot başına \texteuro{} cinsinden mutlak açık olarak karşılaştırılan her yöntem için iyi tanımlıdır:
\begin{equation}\label{eq:gap}
\mathrm{gap}^{\sm{pen}}(\pi)=\frac{J^{\sm{pen,LP}}-J^{\sm{pen}}(\pi)}{N}\ \ge0 .
\end{equation}
Epizot başına göreli bir açık kullanılmamaktadır, çünkü $J^{\sm{pen,LP}}$ çoğu zaman sıfıra yakın ya da negatiftir; ek materyal göreli açıkları yalnızca toplamların oranı olarak raporlamaktadır.
Katı optimum $J^{\sm{LP,ceil}}$, toplam eksik enerjiyi en küçükleyen ilk aşamadan sonra ayrılış gereksinimlerini ve trafo limitini kısıt olarak uygular.
Bu optimum, trafo limitine uyan ve daha fazla eksik enerji bırakmayan her denetleyicinin kârını üstten sınırlar (Önerme~\ref{S-prop:oracle}) ve garantinin maliyeti analizinde kullanılır.

\subsection{Kısmen gözlemlenebilir Markov karar süreci}\label{sec:mdp}

Gelecekteki varışlar, kısa bir pencerenin ötesindeki fiyatlar ve gelecekteki yükler gözlemlenmediğinden çevrim içi problem bir POMDP'dir.
Mevcut gözlem üzerinden ileri beslemeli bir aktör ve kritik ile eylem seçilmektedir.
Politika her şarj ünitesi için bir ham eylem $a_i(t)\in\R$ üretir; bu eylem $P^{\sm{ref}}=11$ kW olmak üzere $x_i(t)=P^{\sm{ref}}a_i(t)$ biçiminde kW'a ölçeklenir ve \Sec{sec:projection} içindeki güvenlik katmanı onu herhangi bir sıkıştırma fonksiyonu kullanmadan komut edilen güce eşler.
Politikanın ödülü adımın gerçekleşen kârıdır; ayrılış gereksinimi ya da trafo limiti için ceza içermez, çünkü güvenlik katmanı her ikisini de uygular.
Eğitimde ödül, getirilerin $N$ ile büyümemesi için $r^{\sm{ref}}=c^{\sm{ref}}P^{\sm{ref}}\Delta t$ ve $c^{\sm{ref}}=0.25$\,\texteuro/kWh olmak üzere $N r^{\sm{ref}}$'ye bölünür; raporlanan tüm sonuçlar ise bölünmemiş kârı kullanır.
Varışlar, fiyatlar, yükler ve PV eylemlerden bağımsız olarak değişir, her EV epizot bitmeden ayrılır ve indirim faktörü $\gamma=0.99$, $\Delta t=15$ dk için yaklaşık bir günlük etkin ufuk sağlar.

% ---------------------------------------------------------------------
\section{Önerilen Yöntem}\label{sec:method}
% ---------------------------------------------------------------------

\Fig{fig:architecture} üzerinde görüldüğü gibi \method{}, her adımda art arda çalışan üç parçadan oluşur.
Dikkat tabanlı politika, bir şarj ünitesi belirteçleri kümesini şarj ünitesi başına bir ham eyleme eşler; güvenlik katmanı ise ham eylemleri araç başına sınırlara ve trafo limitine uyan komut güçlerine dönüştürür.
Eğitim sırasında bir politika gradyanı yöntemi, güvenlik katmanını ortamın bir parçası olarak ele alarak politikayı gerçekleşen kâra göre günceller; aynı katman kullanım sırasında da değişmeden çalışır.

\begin{figure}[t]
\centering
\figorbox{figures/fig_architecture.pdf}{\linewidth}{6cm}
\caption{\method{} yönteminin mimarisi ve kontrol döngüsü. Düz oklar eğitimde ve kullanımda işletilen kontrol döngüsünü oluşturur; kesikli oklar yalnızca eğitimde kullanılır. Güvenlik katmanı (gölgeli), trafo limitini koruyan ve her EV alt sınırını aldığı sürece EV'nin gereksinimini erişilebilir tutan kesin bir izdüşümdür (Teorem~\ref{thm:invariance}).}
\label{fig:architecture}
\end{figure}

% ---------------------------------------------------------------------
\subsection{Belirteç gösterimi ve dikkat tabanlı politika}\label{sec:policy}

Gözlem, $N$ adet şarj ünitesi belirteci $\xi_i(t)\in\R^{12}$ ve bir küresel belirteç $\xi_0(t)\in\R^{38}$ içeren bir kümedir; hiçbir öznitelik bir şarj ünitesinin indisini ya da konumunu kodlamaz.
Bir şarj ünitesi belirteci doluluk, enerji ve kapasite verilerini, kalan gerekli enerjiyi ve adım sayısını, \Sec{sec:bounds} içinde tanımlanan esneklik payını, güç limitlerini ve $l_i$ ile $u_i$ sınırlarını içerir.
Küresel belirteç sonraki $H_f=8$ adımın fiyatlarını, zaman özniteliklerini, mevcut limitleri, yük ve PV tahminlerini ve dört istasyon oranını içerir.
Tüm öznitelikler $N$ ile büyümeyecek biçimde normalize edilir; Tablo~\ref{S-tab:features} bunları listelemektedir.

Her belirteç, öğrenilen iki tür gömmesinden biriyle birlikte doğrusal olarak gömülür: şarj üniteleri için $e_i=W_c\xi_i+b_c+v_c$, küresel belirteç için $e_0=W_g\xi_0+b_g+v_g$.
$[e_0;e_1;\dots;e_N]$ dizisi, genişliği $d_m=128$ olan, $n_H=4$ başlıklı ve konum kodlaması içermeyen $n_L=3$ adet ön-LayerNorm'lu Transformer kodlayıcı katmanından geçer \cite{vaswani2017attention,ba2016layer,xiong2020layer}.
Bir anahtar dolgu maskesi boş şarj ünitelerini dikkat anahtarları arasından çıkarır; küresel anahtar ise hiçbir zaman maskelenmez.
Ardından ortak bir aktör başlığı, her şarj ünitesi çıktısı $h_i$'yi Gauss dağılımlı bir ham eylemin $\mu_{\theta,i}$ ortalamasına ve $\sigma_{\theta,i}$ standart sapmasına eşler; burada $\theta$ aktör parametrelerini gösterir.
Politika $\pi_\theta(a\mid s)=\prod_{i\in\mathcal N(t)}\mathrm{Normal}\big(a_i;\mu_{\theta,i}(s),\sigma_{\theta,i}^2(s)\big)$ biçimindedir.
$\omega$ parametreli ayrı bir kritik kodlayıcısı, maskelenmemiş tüm belirteç çıktılarını öğrenilen tek bir dikkat sorgusuyla \cite{lee2019set} havuzlayarak $V_\omega(s)$ değerini üretir.

Bu yapı permütasyona eşdeğişkendir (permutation-equivariant); maskelenmiş belirteçler maskelenmemiş belirteçlerin çıktılarını etkilemez ve parametre sayısı $N$'ye bağlı değildir (Lemma~\ref{S-lem:equivariance}).
Şarj ünitesi sayısı ne olursa olsun aktörün 610\,370, kritiğin ise 676\,481 parametresi vardır.
Dolayısıyla $N\in\{10,20\}$ ile eğitilen ağırlıklar $N=200$ için çalıştırılabilir; ancak bu boyutta performans garanti edilmez, çünkü 200 belirteç üzerindeki dikkat, eğitimde görülenden daha fazla anahtar üzerinden ortalama alır.

% ---------------------------------------------------------------------
\subsection{Erişilebilir enerjiden türetilen araç başına sınırlar}\label{sec:bounds}

Güvenlik katmanı, bağlı her EV için ayrılış gereksinimini erişilebilir tutan en küçük güce ihtiyaç duyar; karışıklığa yol açmadığı yerlerde $i$ indisi yazılmamaktadır.
Bir adım boyunca tam güçte şarj, $E$ enerjisini $\varphi(E)=E+\Delta t\,\eta^{\sm c}\,\bar P(E)$ değerine eşler.
\Sec{sec:invariance} içindeki A1 varsayımı, $\varphi$'nin sürekli ve azalmayan olmasını, yani bir adıma daha fazla enerjiyle başlamanın tam güçte şarjdan sonra hiçbir zaman daha az enerji bırakmamasını gerektirir.
Gereksinime tam güçte hâlâ erişilebilen en küçük enerjiler olan geriye doğru erişilebilirlik eşikleri şunlardır:
\begin{equation}\label{eq:thresholds}
B(0)=\Ereq,\qquad B(k)=\min\big\{E\in[E^{\min},E^{\sm{cap}}]:\ \varphi(E)\ge B(k-1)\big\},\quad k\ge1 .
\end{equation}
A1 altında bir EV, $t$ adımında bireysel olarak olurludur (yani kabul edilebilir bir kontrol \Eq{eq:requirement} koşulunu sağlar) ancak ve ancak $E(t)\ge B(\tau(t))$ ise (Lemma~\ref{S-lem:reach}).
A1 olmadan bu nitelendirme geçerliliğini yitirir, çünkü şimdi daha az şarj etmek daha sonra daha fazla erişilebilir enerji bırakabilir.
EV2Gym'in üstel şarj gücü azalması için Lemma~\ref{S-lem:taper}, A1'i her adım uzunluğu için ispatlamaktadır.

$E=E_i(t)$ enerjisine ve $\tau=\tau_i(t)$ kalan eyleme sahip bağlı bir EV için güvenlik katmanının sınırları şunlardır:
\begin{equation}\label{eq:bounds}
u_i(t)=\bar P(E),\qquad
l_i(t)=\begin{cases}
\min\Big\{u_i(t),\ \max\Big\{\pmn(E),\ g^{-1}\Big(\dfrac{\Bh(\tau-1)-E}{\Delta t}\Big)\Big\}\Big\}, & E\ge\Bh(\tau),\\[2mm]
\bar P(E), & E<\Bh(\tau),
\end{cases}
\end{equation}
burada $\Bh$ saklanan bir eşik tablosu, dıştaki minimum ise sayısal bir korumadır.
Birinci durumda $l_i$, ardından gelen $\tau-1$ tam güç adımının $\Ereq$ değerine hâlâ ulaştığı kabul edilebilir en küçük güçtür; bireysel olarak olursuz bir EV ise en yüksek güçte şarj edilir.
Negatif bir alt sınır, pay kaldığı sürece V2G deşarjına izin verir ve \Fig{fig:concepts}(a) üzerinde görüldüğü gibi sınır, enerjiyi en geç başlangıç yörüngesinin üzerinde ya da bu yörünge üzerinde tutar.

\begin{figure}[t]
\centering
\figchain{\linewidth}{\analysisdir/figures/fig2_concepts.pdf}{\analysisdir/figures_conceptual/fig_concepts.pdf}{figures/fig2_concepts.pdf}
\caption{Güvenlik katmanının kavramları. (a) Bir EV'nin kalış süresi boyunca enerjisi; en geç başlangıç yörüngesi $\Bh_i$ ve adım başına $\Delta t\,u_i$, $\Delta t\,l_i$ ve $\Delta t\,\pmn_i$ seçenekleriyle birlikte. (b) İki EV için $-\Crev\le p_1+p_2\le\Ctil$ dilimiyle kesilen $[l,u]$ kutusuna izdüşüm.}
\label{fig:concepts}
\end{figure}

$\Bh_i(k)$, $k=0,\dots,\tau_i(t_i^{\sm{arr}})$ tablosu varışta $\varphi_i$ üzerinde ikiye bölme (bisection) yöntemiyle bir kez hesaplanır; böylece sınırların maliyeti EV ve adım başına $O(1)$ olur.
Saklanan her tablo
\begin{equation}\label{eq:certified}
\Bh(0)\ge\Ereq,\qquad \varphi\big(\Bh(k)\big)\ge\Bh(k-1),\qquad \Bh(k)\le\max\{\Bh(k-1),E^{\min}\},\qquad \Bh(k)\in[E^{\min},E^{\sm{cap}}]
\end{equation}
anlamında sertifikalıdır ve ispatlar yalnızca bu özellikleri kullanır.
Adım cinsinden ölçülen esneklik payı $\chi_i(t)=\big(E_i(t)-\Bh_i(\tau_i(t))\big)/\big(\eta_i^{\sm c}\Pnom_i\Delta t\big)$, ancak ve ancak EV bireysel olarak olurlu olduğunda negatif değildir.

% ---------------------------------------------------------------------
\subsection{İzdüşüm}\label{sec:projection}

$x\in\R^N$, kW cinsinden ham eylem olsun; boş şarj üniteleri için $x_i=l_i=u_i=0$ alınır.
Güvenlik katmanı, sıkılaştırılmış ileri yön limiti olan $\Ctil(t)=C(t)-n(t)\varepsilon_q$ ifadesini kullanır; burada şarj ünitesi başına $\varepsilon_q=0.008$ kW'lık pay, simülatörün eyleyici yuvarlamasını karşılar (Bölüm~\ref{S-app:actuation}).
Olurlu küme $\mathcal K(t)=\{p\in\R^N:\ l\le p\le u,\ -\Crev(t)\le\mathbf 1^\top p\le\Ctil(t)\}$ olup güvenlik katmanı şu Öklid izdüşümüdür:
\begin{equation}\label{eq:proj}
\Pi_t(x)=\arg\min_{p\in\mathcal K(t)}\tfrac12\|p-x\|_2^2 .
\end{equation}

\begin{proposition}[İzdüşüm]\label{prop:projection}
\begin{enumerate}
\item[(i)] $\mathcal K(t)$ kümesi ancak ve ancak $l\le u$, $\sum_il_i\le\Ctil$, $\sum_iu_i\ge-\Crev$ ve $-\Crev\le\Ctil$ ise boş değildir.
\item[(ii)] $\mathcal K(t)\neq\emptyset$ ise \Eq{eq:proj} probleminin tek minimize edicisi $p_i^{*}=\clip(x_i-\lambda^{*},l_i,u_i)$'dir. Kırpılmış nokta toplam limiti zaten sağlıyorsa ortak kaydırma $\lambda^{*}$ sıfırdır; aksi hâlde $\lambda^{*}$, $\lambda\mapsto\sum_i\clip(x_i-\lambda,l_i,u_i)$ fonksiyonunun ihlal edilen limitteki köküdür.
\end{enumerate}
\end{proposition}

İspat Karush--Kuhn--Tucker (KKT) koşullarını doğrulamaktadır; ek materyal ayrıca Jacobi matrisini de vermektedir (Önerme~\ref{S-prop:jacobian}).
\Fig{fig:concepts}(b) üzerinde görüldüğü gibi izdüşüm, tüm serbest koordinatları aynı miktarda kaydırır ve toplam limit sağlanana kadar kırpar.
$2N$ kırılma noktasının $\{x_i-u_i\}\cup\{x_i-l_i\}$ bir kez sıralanması ve bunlar üzerinden bir kez geçilmesi $\lambda^{*}$ değerini verir; dolayısıyla katmanın maliyeti adım başına $O(N\log N)$'dir.
İkiye bölme yöntemine dayalı yedekli bir son denetim, kayan nokta yuvarlamasına karşın toplam limiti korur; Algoritma~\ref{S-alg:safety} tam bir adımı listelemektedir.

Bu cebir yeni değildir: \Eq{eq:proj} sürekli kuadratik sırt çantası problemidir \cite{kiwiel2007breakpoint,wang2022efficient}; sLLF hız kuralı da buna çok yakın bir biçim, yani her EV için tepe hızıyla ölçeklenen ortak bir düzey kullanır \cite{chen2022smoothed}.
Bizim kullanımımız iki bakımdan farklıdır: girdi, boyuttan bağımsız bir politikanın öğrenilmiş şarj ünitesi başına eylemidir; $l_i$ alt sınırı ise şarj gücü azalmasını ve V2G işaret aralığını içeren geriye doğru erişilebilirlikten türetilir.
Ayrıca toplam kısıtı iki yönlüdür ve katman eğitim sırasında da çalışır.

% ---------------------------------------------------------------------
\subsection{Olurluluğun korunması}\label{sec:invariance}

Sezgisel olarak güvenlik katmanı, hiçbir EV'nin gereksinimine tam güçte hâlâ erişilebilen en küçük enerjinin altına düşmesine izin vermez; ortak kaydırma ise istasyonu trafo bandı içinde tutar.
Bu özelliklerden yalnızca ilki sağlanamayabilir; bu durum, trafo her EV'ye alt sınırını aynı anda veremediğinde ortaya çıkar.
$\mathcal K(t)$ boş olmadığında bir adımı olurlu olarak adlandırıyoruz; olursuz bir adımda izdüşümün yerini bir yedek kural alır.
Sonucun koşulları aşağıda toplanmıştır ve bunlardan yalnızca ilk beşi varsayımdır.
Teknik açıdan, sonucun dayandığı tümevarım Chen ve arkadaşlarının esneklik payı argümanını \cite{chen2022smoothed} şarj gücü azalmasına, V2G deşarjına, iki yönlü bir bağlantıya ve keyfî ham eylemlere genişletmektedir.

\begin{center}
\textbf{Teorem~\ref{thm:invariance}'in koşulları}\\[1mm]
\footnotesize
\setlength{\tabcolsep}{3pt}
\begin{tabular}{@{}l>{\raggedright}p{6.4cm}>{\raggedright\arraybackslash}p{6.2cm}@{}}
\toprule
Etiket & Koşul & Bu çalışmadaki durumu \\
\midrule
A1 & $\varphi$ sürekli ve azalmayan & EV2Gym'in şarj gücü azalması için ispatlanmıştır (Lemma~\ref{S-lem:taper}) \\
A2 & ayrılış adımı $d_i$ ve gereksinim $\Ereq_i$ varışta biliniyor & konformal sınır ile gevşetilmiştir (\Sec{sec:conformal}) \\
A3 & tek adım modeli: \Eq{eq:limits} içindeki komutlar için $E_i(t+1)\ge E_i(t)+\Delta t\,g_i(p_i)$ ve $p_i\le\tilde p_i\le p_i+\varepsilon_q$ & EV2Gym'in eyleyici davranışı için ispatlanmıştır (Lemma~\ref{S-lem:a3}); donanımda sağlanamayabilir \\
A4 & olurlu varış: $E_i^{\sm{arr}}\ge\Bh_i(\tau_i(t_i^{\sm{arr}}))$ & talep işleme hattı tarafından uygulanır \\
A6 & şebeke kapasite payı: $C(t)\ge N\varepsilon_q$ ve $\Crev(t)\ge0$ & dışsaldır; yedek kural sayımlarıyla denetlenir \\
\midrule
A5 & olurlu adım: $\sum_il_i\le\Ctil$, $\sum_iu_i\ge-\Crev$, $\Ctil\ge-\Crev$ & kapalı döngüde ortaya çıkan bir olaydır; (a) ve (b) tarafından varsayılmaz; ölçülür \\
\bottomrule
\end{tabular}
\end{center}

A6, politikaya bağlı olmaması için $n(t)$ yerine şarj ünitesi sayısı $N$'yi kullanır ve senaryolarımızda epizot başına $2\times10^{-9}$'un altında bir olasılıkla sağlanamaz (Bölüm~\ref{S-app:a6}).

\begin{theorem}[Olurluluğun korunması]\label{thm:invariance}
Eğitim ya da kullanım sırasında, sonlu ham eylemlerden oluşan keyfî bir dizi için komut edilen güçler olurlu adımlarda $p(t)=\Pi_t(x(t))$, diğer adımlarda ise yedek kuralın çıktısı olsun.
\begin{enumerate}
\item[(a)] A3 ve A6 altında, her adımda gerçekleşen güçler $-\Crev(t)\le\sum_i\tilde p_i(t)\le C(t)$ trafo limitini, komut edilen güçler ise \Eq{eq:limits} ile verilen fiziksel limitleri sağlar.
\item[(b)] A1--A4 altında, kalış süresi boyunca $p_i(t)\ge l_i(t)$ olan her EV $i$, $E_i(d_i+1)\ge\Ereq_i$ koşulunu sağlar; özellikle kalış süresi boyunca olursuz adımla karşılaşmayan her EV bu koşulu sağlar.
\item[(c)] Ek olarak her adım olurlu ise her EV ayrılış gereksinimini karşılar.
\end{enumerate}
\end{theorem}

\begin{proof}[İspat taslağı]
Olurlu adımlarda Önerme~\ref{prop:projection}(i) $p(t)\in\mathcal K(t)$ sonucunu verir ve A3, sıkılaştırılmış bandı gerçekleşen güçlere aktarır.
A6 altında yedek kural yalnızca aşağıda açıklanan gerekli güç fazlası durumuyla karşılaşabilir ve bu durumda çıktısı, komut edilen toplamı $\Ctil(t)$ değerinde tutar.
(b) için tek adımlık bir lemma (Lemma~\ref{S-lem:onestep}), $E_i(t)\ge\Bh_i(\tau_i(t))$ değişmezini A4'ün sağladığı varıştan itibaren $p_i\ge l_i$ olan her adım boyunca taşır; son adımda bu değişmez $E_i(d_i+1)\ge\Bh_i(0)\ge\Ereq_i$ sonucunu verir.
(c) şıkkı (b)'den çıkar; tam ispat Bölüm~\ref{S-app:theorem} içinde verilmektedir.
\end{proof}

Trafoya ilişkin kısım ne politikaya ne de A1 ya da A5'e bağlıdır; ayrılışa ilişkin kısım ise araç başına bir ifadedir.
A5, ortak bir plan mevcut olduğunda bile sağlanamayabilir.
10 kW'lık bir limit altında her biri 1 saatlik iki adım içinde 10 kWh'e ihtiyaç duyan iki EV, kapasiteyi eşit biçimde paylaşabilir.
Ancak birinci adımda boşta bekleyen bir politika, ikinci adımda $\sum_il_i=20$ kW bırakır.
Bu nedenle EV'lerin olursuz adımlarla ne sıklıkta karşılaştığı ölçülmektedir.

% ---------------------------------------------------------------------
\subsection{Yedek kural}\label{sec:fallback}

Olursuz bir adımda yedek kural, izdüşümün yerine en az esneklik payı önce ilkesine dayanan deterministik bir paylaştırma koyar.
İlgili durumda gerekli güç $\sum_il_i$, $\Ctil$ değerini aşar.
Bu durumda her EV fiziksel alt limiti $\pmn_i$'den başlar ve $\Ctil$ değerine kadar kalan bütçe, esneklik payı artan sırada EV'lere, her biri en fazla $l_i$'ye kadar olmak üzere dağıtılır.
Bu çıktı trafo limitini korur ve hâlâ $l_i$ alan her olurlu EV olurlu kalır (Lemma~\ref{S-lem:fallback}).
Üç ek durum çelişkili limitleri, aşırı PV ihracatını ve negatif bir bütçeyi kapsar; A6 bu üçünü de dışlar (Bölüm~\ref{S-app:fallback}).
$l_i$'den daha az alan EV'ler için bir epizottaki toplam ayrılış eksikliği, bu adımlarda biriken açık ile önsel olarak sınırlanır (Lemma~\ref{S-lem:deficit}).
Talep işleme hattındaki kabul kontrolü olursuz adımları daha seyrek hâle getirir; ancak tahminlere dayanır ve Teorem~\ref{thm:invariance} kapsamına girmez.

% ---------------------------------------------------------------------
\subsection{Ayrılış adımı için konformal alt sınır}\label{sec:conformal}

Ayrılış adımı açısından A2'yi gevşetmek için tek yönlü, kantil regresyonu biçiminde \cite{romano2019conformalized} bölünmüş konformal tahmin (CP) \cite{papadopoulos2002inductive,vovk2022algorithmic,lei2018distributionfree,angelopoulos2023conformal} kullanılmaktadır.
Bir gradyan artırma (gradient boosting) regresörü $\hat q$, $d_i$'nin alt kantilini varış adımından, haftanın gününden, varış enerjisinden, talep edilen enerjiden ve EV modelinden kestirir.
Kalibrasyon, bağımsız tohumlarla (seed) üretilmiş $n_{\sm{cal}}$ epizodun her birinden düzgün dağılımla bir oturum seçer ve $e_j=\hat q(\mathbf w_j)-d_j$ skorlarını kullanır.
$\hat Q_{1-\alpha}$, $\lceil(n_{\sm{cal}}+1)(1-\alpha)\rceil$'inci en küçük skor olmak üzere $\hat d_i^{\sm{lo}}=\max\{t_i^{\sm{arr}},\lceil\hat q(\mathbf w_i)-\hat Q_{1-\alpha}\rceil\}$ sınırı; sınırlarda, talep işleme hattında ve özniteliklerde $d_i$'nin yerini alır ve $\hat d_i^{\sm{lo}}$ adımından sonra katman tek bir kalan adım kullanır.

\begin{proposition}[Konformal ayrılış garantisi]\label{prop:conformal}
Kalibrasyon oturumları ile test oturumu değiştirilebilir (exchangeable) olsun.
Bu durumda $\mathbb P(d_i\ge\hat d_i^{\sm{lo}})\ge1-\alpha$ olur ve $d_i$ yerine $\hat d_i^{\sm{lo}}$ alınarak A1, A3 ve A4 altında $\mathbb P\big(E_i(d_i+1)\ge\Ereq_i\big)\ge1-\alpha-\mathbb P\big(\exists\,t\le d_i:\ p_i(t)<l_i(t)\big)$ sağlanır.
\end{proposition}

Burada $\Ereq_i$, $\hat d_i^{\sm{lo}}$ ile yeniden sertifikalanmış gereksinimdir ve olasılık, kalibrasyon verisi ile bir test epizodundan düzgün dağılımla seçilen bir EV üzerinden alınır; dolayısıyla garanti marjinaldir.
Kalibrasyon her test yapılandırmasından veri gerektirdiğinden konformal varyant sıfır-atış değildir.

% ---------------------------------------------------------------------
\subsection{Eğitim algoritması}\label{sec:training}

Politika, güvenlik katmanı, yedek kural ve eyleyici davranışı ortamın parçası olarak ele alınarak yakınsal politika optimizasyonu (PPO) \cite{schulman2017proximal} ve genelleştirilmiş avantaj kestirimi \cite{schulman2016highdimensional} ile eğitilir.
Bu nedenle PPO, ham Gauss örneğinin olabilirlik oranını kullanır.
Eski parametrelerdeki PPO vekil amaç gradyanına eşit olan temel (vanilla) politika gradyanı kestiricisi, bu durumda izdüşümü alınmış eylemlerin getirisi için yansızdır.
Stokastik politikalar için ortamın içindeki ve politikanın içindeki izdüşüm bu anlamda denktir \cite{markgraf2025safe}; bu tercih aktif bir izdüşümün sıfır gradyanlarından da kaçınır \cite{kasaura2023benchmarking}.
$u_i-l_i\le10^{-9}$ kW olan belirteçler hiçbir komut gücünü etkileyemeyeceğinden log-olasılıktan ve entropiden çıkarılır.
Yedek kural adımları aktör kaybında maskelenir, ancak kritik için tutulur.
Birçok ham eylem aynı izdüşüm noktasına eşlendiğinden olurlu adımlarda $\beta=0.01$ ağırlıklı bir düzenlileştirici, politika ortalamasını $\mathcal K(t)$'ye yakın tutar:
\begin{equation}\label{eq:reg}
\mathcal L_{\sm{reg}}(\theta)=\beta\,\mathbb E_{s}\Big[\frac1{n(t)}\sum_{i\in\mathcal N(t)}\Big(\frac{\bar x_i-\Pi_t(\bar x)_i}{P^{\sm{ref}}}\Big)^{2}\Big],\qquad \bar x=P^{\sm{ref}}\mu_\theta(s).
\end{equation}
Moreau özdeşliğine göre $\tfrac12\mathrm{dist}^2(x,\mathcal K)$'nin gradyanı $x-\Pi_t(x)$'tir; bu nedenle $\Pi_t(\bar x)$ gradyan durdurma (stop-gradient) altında hesaplanır \cite{chen2021enforcing,markgraf2025safe}.
Her veri toplama turu (rollout) $N=10$ olan sekiz ve $N=20$ olan sekiz ortam çalıştırır; tam eğitim döngüsü Algoritma~\ref{S-alg:ppo} içinde verilmektedir.

% ---------------------------------------------------------------------
\section{Deney Düzeneği}\label{sec:setup}
% ---------------------------------------------------------------------

\subsection{Simülatör ve senaryolar}\label{sec:scenarios}

Tüm deneylerde, kaynak kodunda değişiklik yapılmadan \texttt{55ef368c} commit'indeki EV2Gym \cite{orfanoudakis2025ev2gym} kullanılmaktadır.
Bir sarmalayıcı (wrapper) gözlemleri hesaplar; ortak kabul edilen gereksinimleri, güvenlik katmanını, yedek kuralı ve eyleyici davranışını uygular ve gelecekteki varışları okumadan tüm ölçütleri kaydeder.
Her istasyonda tek bir trafonun arkasında $N$ şarj ünitesi bulunur; EV'ler simülatörün V2G destekli sekiz EV modelinden, her model için sabit bir verimle (bir model için 0,93, diğerleri için 0,90) çekilir.
İstenen enerji bataryanın tam dolu hâlidir; fiyatlar EV2Gym ile birlikte sunulan Hollanda gün öncesi fiyatlarıdır ve deşarj fiyatı şarj fiyatına eşittir.
Tablo~\ref{S-tab:scenarios} tüm senaryo parametrelerini listelemektedir.

Şebeke büyüklükleri ve varış hızı $N$ ile doğrusal olarak ölçeklenir.
Özellikle trafo anma gücü $P_{\sm{tr}}=\rho_{\sm{tr}}N$'dir; burada EV2Gym'in 25 şarj üniteli V2G yapılandırmasındaki gibi şarj ünitesi başına $\rho_{\sm{tr}}=4$ kW'lık sıkı bir düzey ve 7 kW'lık orta bir düzey kullanılır.
Esnek olmayan yük ve PV, ortalaması $0.4P_{\sm{tr}}$ olan çarpanlarla ölçeklenir ve gözlemdeki tahminler standart sapması \%10 olan Gauss gürültüsü içerir.
EV2Gym'in işyeri, kamusal ve özel varış profilleri, gece boyu süren oturumların 15 dakikalık $T=112$ adımdan oluşan 28 saatlik ufka sığmasını sağlayan saatlerde başlar.

Eğitim, her yığında (batch) $N\in\{10,20\}$ değerlerini, üç profili ve iki anma düzeyini karıştırır.
Test ızgarası $N\in\{5,10,20,50,100,200\}$ değerlerini, her iki anma düzeyini ve tüm profilleri; buna ek olarak $N\in\{20,200\}$ için $\rho_{\sm{tr}}=3$ kW'lık bir stres düzeyini ve EV2Gym'in özgün 25 şarj üniteli yapılandırmasını, toplamda 43 yapılandırmayı kapsar.
Dolayısıyla sıfır-atış ifadesi istasyon boyutunu ve stres düzeyini ifade eder; profiller ise eğitimde görülmüştür.
Her yapılandırma, tüm yöntemlerce paylaşılan ve $N\le50$ için 100, $N\ge100$ için 50 sabit test epizoduna sahiptir; böylece her politika 3.550 test epizodunda çalıştırılır.
Stres düzeyi, çalıştırılan tek duyarlılık çalışmasıdır.

\subsection{Karşılaştırılan yöntemler}\label{sec:baselines}

Her yöntem ve her değerlendirme geçişi aynı kabul edilen gereksinimleri alır ve eksik enerji her zaman bu gereksinimlere göre ölçülür.
Katman içermeyen yöntemler, $l_i$ ya da trafo bağlantısı olmaksızın yalnızca fiziksel kırpma $\clip(x_i,\pmn_i,u_i)$ işlemini uygular.

Öğrenen yöntemler arasında \method{} ve Deep Sets+safety layer, güvenlik katmanına sahip iki boyuttan bağımsız küme politikasıdır.
Attention+penalty, \method{} yönteminin cezaya dayalı ikizidir: güvenlik katmanı olmaksızın, kârdan eksik enerji ve aşırı yük enerjisinin $w$ ağırlığıyla çarpımı çıkarılarak elde edilen değerden öğrenir.
Bunun kısıtlı karşılığı olan Attention+PPO-Lagrangian, sabit ağırlığın yerine maliyet limiti sıfır olan Lagrange çarpanlarını kullanır \cite{altman2021constrained,tessler2018reward}.
$l_i=\pmn_i$ alındığında ``Attention+safety layer without departure bound'' (ayrılış sınırı olmayan dikkat+güvenlik katmanı) varyantı, katmanı \cite{wan2025logicinformed} çalışmasındaki gibi adım başına bir trafo izdüşümüne indirger ve ayrılış gereksinimini ceza içeren ödülden öğrenir.
Ceza ile eğitilmiş kontrol noktaları ayrıca güvenlik katmanı yalnızca kullanım aşamasında uygulanarak çalıştırılır; \method{} ve Attention+penalty ayrıca kabul kontrolü olmaksızın, varışta sınırlandırılmış talep korunarak çalıştırılır.
20 yuvaya kadar sıfırla doldurulmuş belirteçler üzerinde çalışan, katmanlı ve katmansız çok katmanlı algılayıcı (MLP), yalnızca $N\le20$ için değerlendirilen sabit boyutlu bir referanstır.
Dağılım içi bir referans olarak \method{} ayrıca $N=50$ için eğitilir ve $N\in\{50,100,200\}$ için değerlendirilir.

Karşılaştırmayı, her biri katmanlı ve katmansız olmak üzere altı kural tamamlar: olabildiğince hızlı (AFAP), olabildiğince geç (ALAP), döngüsel sıralama (round robin), LLF, sLLF \cite{chen2022smoothed} ve bir fiyat eşiği kuralı.
EV2Gym'in ekonomik MPC'sinin \cite{diazlondono2026open} HiGHS ile yaptığımız yeniden gerçeklemesi eMPC-HiGHS, her $N$ için 10 adımlık bir ufukla çalışır ve bu ufuk içindeki kesin fiyatları ve varışları bilir.
Konformal varyant \method{}-CP, $\alpha=0.1$ için bir kez eğitilir ve her $\alpha\in\{0.01,0.05,0.1,0.2\}$ için yeniden kalibre edilmiş sınırlarla değerlendirilir; ek bir geçiş olarak kalibre edilmemiş bir nokta tahmini de kullanılır.
EV-GNN \cite{orfanoudakis2025scalable} harici bir karşılaştırma yöntemi olarak planlanmış ancak çalıştırılmamıştır; GNN-DT \cite{orfanoudakis2026graph} ise her senaryo için bir optimizasyon çözücüsünün ürettiği veri kümeleri üzerinde çevrim dışı eğitildiğinden uyarlanmamıştır.
Planlanan ablasyonlar, ceza ağırlığı duyarlılıkları, diğer MPC varyantları ve diğer duyarlılık çalışmaları da çalıştırılmamıştır.

\subsection{Ölçütler}\label{sec:metrics}

Ana ölçüt, her yöntem için geçerli olan \Eq{eq:gap} ile tanımlı mutlak açıktır.
Cezalı kârlar, her yöntem ve geçiş için ortak kabul edilen gereksinime göre karşılanmayan enerjiyi kullanır; kâr, karşılanmayan enerji ve trafo aşırı yük enerjisi ayrıca raporlanır.
Güvenlik katmanı için kalış süresi boyunca olursuz adımla karşılaşmayan EV'lerin payı (E4), yedek kural durumlarına göre olursuz adımlar ve Lemma~\ref{S-lem:deficit} ile verilen açık sınırıyla karşılaştırmalı olarak ölçülen eksik enerji raporlanır.
Kayıtlar Teorem~\ref{thm:invariance}(b) hipotezi için EV başına bir bayrak taşımadığından, teoremin ayrılışa ilişkin kısmı epizot düzeyinde denetlenmektedir.

Garantinin maliyeti, aynı epizotlar üzerinde cezaya dayalı ikiz ile güvenlik katmanlı politika arasındaki farktır; birincil ölçü olarak $J^{\sm{pen}}$, ikincil ölçü olarak kâr kullanılır.
Bu maliyet, eşleştirilmiş epizotlar üzerinden toplamların oranıdır ve kabul edilen EV başına \texteuro{} ile $|J^{\sm{LP,ceil}}|$'nin yüzdesi olarak ifade edilir.
Pozitif bir değer, cezaya dayalı ikizin daha yüksek puan aldığı anlamına gelir.
\method{} yönteminin $N=50$ için eğitilen referansa göre aktarım açığı, aktarım kaybını yöntem kaybından ayırır.

\subsection{Ayarlama, istatistik ve hesaplama}\label{sec:hyper}

Ayarlanan yedi öğrenen yöntemin her biri, öğrenme oranı, mini yığın boyutu ve entropi katsayısı üzerinde tanımlı bir ızgaradan 12 yapılandırma almıştır.
Ceza içeren ödüller ızgaraya $w\in\{1,10,100\}$ ağırlığını, Lagrange yöntemi ise çarpan adım büyüklüğünü ekler.
Her yapılandırma eğitim karışımı üzerinde $0.5\times10^{6}$ adım eğitilmiş ve test tekrarlarından ayrık 24 doğrulama epizodunda değerlendirilmiştir; seçimde her yöntem için $w=10$ ile hesaplanan doğrulama $J^{\sm{pen}}$ değeri kullanılmıştır.
Seçilen yapılandırma iki tohumla $3\times10^{6}$ adım boyunca yeniden eğitilmiştir; Tablo~\ref{S-tab:hyper} hiperparametreleri listelemektedir.

İki tohumla tohumlar arası değişkenlik test edilemediğinden analiz, tohum ortalamalı politikaları ortak test tekrarları üzerinde karşılaştırmaktadır.
Tablolar, tohum ortalamalı politikanın ortalamasını iki tohum üzerindeki aralıkla birlikte raporlamaktadır.
Tekrar düzeyindeki farklar üzerinde eşleştirilmiş Wilcoxon işaretli sıralar testleri \cite{wilcoxon1945individual}, her istasyon boyutu için \method{} yöntemini diğer her yöntemle karşılaştırır.
Holm düzeltmesi bir sonuç tablosunun tüm karşılaştırmaları üzerinden (aile başına en fazla 351) uygulanır ve sıra-çift serili korelasyonun (rank-biserial) işareti yönü verir.
Ortalamanın güven aralıkları, test tekrarlarını her yapılandırma içinde yeniden örnekleyen bir bootstrap ile elde edilir ve açığın çeyrekler arası ortalaması (IQM) eğitim boyutlarında raporlanır \cite{agarwal2021deep}.
Wilcoxon testi medyan kaymasıyla, güven aralığı ise ortalamayla ilgilidir; büyük cezalı birkaç epizot tek başına ortalamayı kaydırabilir.
İkisinin uyuşmadığı durumlarda her ikisi de raporlanmakta ve karşılaştırma sonuçsuz (inconclusive) olarak nitelendirilmektedir.
Az sayıda tohuma dayanan sonuçlar yanıltıcı olabilir \cite{henderson2018deep}; bu nedenle tohum sayısı bir sınırlılıktır (Bölüm~\ref{S-app:stats}).

Gerçekleme PyTorch \cite{paszke2019pytorch}, CleanRL tarzında \cite{huang2022cleanrl} tek dosyalık bir PPO ve Adam optimize edicisini \cite{kingma2015adam} kullanmakta; güvenlik katmanı çift duyarlıkla çalışmaktadır.
Tüm deneyler WSL2 Ubuntu 24.04 altında Python 3.13.15 ve PyTorch 2.14.1 ile bir AMD Ryzen 7 7700 işlemci ve bir NVIDIA GeForce RTX 5070 Ti (16 GB) üzerinde çalıştırılmıştır.
Hazırlık 1,6 saat, ayarlama 7,9 saat, eğitim 10,2 saat ve değerlendirme 2,5 saat sürmüş; toplamda 153.100 epizot çalıştırılmıştır (Bölüm~\ref{S-app:compute}).
Kod, yapılandırmalar ve analiz betikleri makaleyle birlikte yayımlanacaktır \TODO{depo bağlantısı}.

% ---------------------------------------------------------------------
\section{Bulgular}\label{sec:results}
% ---------------------------------------------------------------------
% İSKELET. Henüz hiçbir deney çalıştırılmamıştır. Her nicel ifade, yalnızca
% analysis/scripts/make_all.py tarafından üretilen değerleri isteyen bir
% \RES{} yer tutucusudur (T1-T7, T3b, TS1-TS3 tabloları; fig3-fig6,
% figS1-figS6 şekilleri; stats/tests_T*.csv). Üretilen tablolar ve şekiller,
% \analysisdir altında (varsayılan ../../analysis) ya da tables/ ve figures/
% içindeki yerel kopyalar olarak var olduklarında otomatik olarak görünür.

\TODO{İskelet: yazarlar aşağıdaki her cümle kalıbını kendi metinleriyle değiştirir, kırmızı yer tutucuları analiz çıktılarından doldurur ve tüm yorumları kendileri yazar}.

\subsection{Ana karşılaştırma}\label{sec:res-main}

\Tab{tab:t1}, profiller ve anma düzeyleri birleştirilerek eğitim boyutları $N\in\{10,20\}$ için tüm yöntemleri karşılaştırmaktadır.
Bu boyutlarda \method{} yönteminin açığı \RES{açık ve \%95 güven aralığıyla IQM, \Tab{tab:t1}} olarak bulunmuştur.
Diğer öğrenen yöntemlere göre cezalı kârdaki fark \RES{farklar, güven aralıkları ve Holm düzeltmeli $p$-değerleri, stats/tests\_T1.csv} şeklindedir.
Kurallar ve model öngörülü yöntemler arasında \RES{güvenlik katmanıyla en iyi kural ve eMPC-HiGHS açıkları, \Tab{tab:t1}}.

\restable{T1}{%
\begin{table}[t]
\centering
\caption{Eğitim istasyon boyutları $N\in\{10,20\}$ için ana karşılaştırma (üç varış profili ve her iki sıkılık düzeyi birleştirilmiştir; analiz betikleri çalıştırılana kadar yer tutucudur).}
\label{tab:t1}
\footnotesize
\setlength{\tabcolsep}{3pt}
\begin{tabular}{@{}lcccccc@{}}
\toprule
Yöntem & $J^{\sm{pen}}$ [\texteuro] & Açık [\%] & IQM açık [\%] & Kâr [\texteuro] & Karşılanmayan [kWh] & Aşırı yük [kWh] \\
\midrule
\method{} (önerilen) & \RES{} & \RES{} & \RES{} & \RES{} & \RES{} & \RES{} \\
Attention+penalty & \RES{} & \RES{} & \RES{} & \RES{} & \RES{} & \RES{} \\
Attention+PPO-Lagrangian & \RES{} & \RES{} & \RES{} & \RES{} & \RES{} & \RES{} \\
Deep Sets+safety layer & \RES{} & \RES{} & \RES{} & \RES{} & \RES{} & \RES{} \\
Kurallar ve eMPC-HiGHS & \RES{} & \RES{} & \RES{} & \RES{} & \RES{} & \RES{} \\
\bottomrule
\end{tabular}
\end{table}}

\subsection{İstasyon boyutları arasında sıfır-atış aktarım}\label{sec:res-scaling}

\Fig{fig:zeroshot}, $N\in\{10,20\}$ için eğitilmiş ağırlıklarla şarj ünitesi başına cezalı kârı ve açığı $N$'ye göre çizmektedir.
$N=50$, $100$ ve $200$ için \method{} yönteminin açığı \RES{\%95 güven aralığıyla her $N$ için açık, \Tab{S-tab:t2}}, $N=50$ için eğitilen referansa göre aktarım açığı ise \RES{\Tab{S-tab:t2}} olarak bulunmuştur.
EV-GNN \RES{kendi eğitim yapılandırmasında ve sıfır-atış durumunda açık, \Tab{S-tab:t2}} değerine ulaşmakta; $\rho_{\sm{tr}}=3$ kW stres düzeyinde \method{} yönteminin açığı ise \RES{\Tab{S-tab:ts2}} olmaktadır.

\begin{figure}[t]
\centering
\resultfig{fig3_zero_shot_scaling.pdf}{Yöntem başına $N$'ye göre şarj ünitesi başına cezalı kâr ve optimallik açığı, \%95 güven aralıkları; eğitim boyutları işaretlenmiştir.}
\caption{5 ile 200 arasındaki $N$ değerleri için istasyon boyutları arasında sıfır-atış aktarım.}
\label{fig:zeroshot}
\end{figure}

\subsection{Kısıtların sağlanması}\label{sec:res-constraints}

\Fig{fig:constraints}, ana yöntemler için karşılanmayan enerjiyi, aşırı yük enerjisini ve olursuz adımla karşılaşmayan EV'lerin payını (E4) $N$'ye göre göstermektedir.
Güvenlik katmanıyla aşırı yük enerjisi \RES{\Tab{S-tab:t3} içindeki değer; Teorem~\ref{thm:invariance}(a) sıfır öngörür} olmakta; katman olmadan ise yöntemler \RES{karşılanmayan enerji ve aşırı yük, \Tab{S-tab:t3}} bırakmaktadır.
\method{} yönteminin E4 değeri \RES{her $N$ için, \Tab{S-tab:t3b}} olup toplam limitler \RES{her $N$ için bağlayıcılık oranları, \Tab{S-tab:t3b}} oranında bağlayıcı olmaktadır.
Olursuz adımlar yedek kural durumlarına \RES{durum başına sayılar, \Tab{S-tab:t3b}} biçiminde dağılmakta ve eksik enerji \RES{açık sınırına oran, \Tab{S-tab:t3}} düzeyine ulaşmaktadır.

\begin{figure}[t]
\centering
\resultfig{fig4_constraint_satisfaction.pdf}{Ana yöntemler için $N$'ye göre karşılanmayan enerji, aşırı yük enerjisi ve E4.}
\caption{Kalış süresi boyunca olursuz adımla karşılaşmayan EV'lerin payı (E4) dâhil olmak üzere istasyon boyutları arasında kısıtların sağlanması.}
\label{fig:constraints}
\end{figure}

\subsection{Garantinin maliyeti}\label{sec:res-cost}

\Fig{fig:cost}, \Sec{sec:metrics} içinde tanımlanan garantinin maliyetini $N$'ye göre çizmektedir.
Cezalı kâr açısından Attention+penalty eksi \method{} farkı \RES{her $N$ için hizmet verilen EV başına \texteuro{} ve $|J^{\sm{LP,ceil}}|$'nin \%'si olarak $\Delta J^{\sm{pen}}$, \Tab{S-tab:t5}} olarak bulunmuştur.
Ham kâr cinsinden ölçüldüğünde fark \RES{her $N$ için $\Delta$ kâr, \Tab{S-tab:t5}} olup testler \RES{stats/tests\_T5.csv} içinde verilmektedir.

\begin{figure}[t]
\centering
\resultfig{fig6_cost_of_guarantee.pdf}{$N$'ye göre garantinin maliyeti ($\Delta J^{\sm{pen}}$ ve $\Delta$ kâr, hizmet verilen EV başına \texteuro{}).}
\caption{İstasyon boyutları arasında garantinin maliyeti. Pozitif bir değer, aynı epizotlar üzerinde Attention+penalty'nin \method{} yönteminden daha yüksek puan aldığı anlamına gelir.}
\label{fig:cost}
\end{figure}

\subsection{Diğer bulgular}\label{sec:res-further}

\Tab{S-tab:t4} içinde ayrılış sınırı olmayan, katmanı yalnızca kullanım aşamasında uygulayan ve kabul kontrolü içermeyen varyantlar \RES{\method{} yöntemine göre açıklar, \Tab{S-tab:t4}} sonucunu vermektedir.
Konformal varyant için açık ve epizot ortalamalı kapsama \RES{her $\alpha$ için, \Tab{S-tab:t7}} olup $N=200$ için bir kontrol adımı \RES{çıkarım süresi, \Tab{S-tab:t6}} sürmektedir.
Duyarlılık çalışmaları \RES{\Tab{S-tab:ts2}} sonucunu vermekte; Şekil~\ref{S-fig:exampleday} ise örnek bir epizodu göstermektedir.

% ---------------------------------------------------------------------
\section{Tartışma}\label{sec:discussion}
% ---------------------------------------------------------------------

\subsection{Yorum}\label{sec:interpretation}

% Bulguların yorumu yazarlar tarafından yazılmalıdır (SAUCIS yapay zekâ
% politikası). Yönlendirmeler, her paragrafın dayandığı analiz çıktısını
% belirtir; burada önceden yazılmış hiçbir yorum yoktur.

\TODO{Paragraf 1, yazarlar tarafından: 10 ve 20 şarj ünitesiyle eğitilen ağırlıkların 200 şarj ünitesine ve stres düzeyine aktarılıp aktarılmadığını belirtin. \Fig{fig:zeroshot}, \Tab{S-tab:t2} içindeki aktarım açığı ve yedek kural sütunları, \Tab{S-tab:t3b} içindeki bağlayıcılık oranları ve \Tab{S-tab:ts2} kullanılmalıdır}.

\TODO{Paragraf 2, yazarlar tarafından: \Fig{fig:cost} ve \Tab{S-tab:t5} içindeki garantinin maliyetini, \Tab{S-tab:t3} içindeki Attention+penalty ihlalleriyle ilişkilendirin. Birincil ölçü olarak $\Delta J^{\sm{pen}}$ ve stats/tests\_T5.csv kullanılmalıdır}.

\TODO{Paragraf 3, yazarlar tarafından: \Tab{S-tab:t3} ve \Tab{S-tab:t3b} içindeki E4'ü ve açık sınırına göre eksik enerjiyi tartışın. Her ikisini A5'in kapalı döngü niteliğiyle ilişkilendirin}.

\TODO{Paragraf 4, yazarlar tarafından: yayımlanmış EV2Gym sayılarıyla karşılaştırmayı yalnızca literatür yapılandırması üzerinde (\Tab{S-tab:ts3}) yapın; farklı simülatörler ya da senaryolar arasında asla karşılaştırma yapmayın}.

\subsection{Şarj noktası işletmecileri için pratik çıkarımlar}\label{sec:implications}

Bu alt bölüm, deneylerden sonra yazarlar tarafından aşağıdaki noktaları kapsayacak şekilde yazılacaktır.
\begin{itemize}
\item \TODO{Trafo limiti her ham eylem için sağlandığında bir CPO'nun ne kazandığı ve koşullu ayrılış sonucunun hizmet taahhüdü açısından ne anlama geldiği.}
\item \TODO{\Tab{S-tab:t5} üzerinden, belirli bir boyuttaki istasyon için hizmet verilen EV başına avro cinsinden garantinin maliyeti.}
\item \TODO{\Tab{S-tab:t7} üzerinden, sürücülerin ayrılış zamanlarını bildirmediği durumlarda konformal varyantın kullanımı ve kabul edilebilir kapsama düzeyleri.}
\item \TODO{\Tab{S-tab:t6} üzerinden, 15 dakikalık kontrol adımına göre kontrol adımı başına çıkarım süresi.}
\end{itemize}

\subsection{Sınırlılıklar}\label{sec:limitations}

Teorem~\ref{thm:invariance} kapsamındaki ayrılış sonucu, bildirilen ayrılış zamanlarına ve gereksinimlere (A2) dayanmaktadır.
Konformal genişletmemiz yalnızca ayrılış adımını gevşetmekte, kapsaması marjinal kalmakta ve son tarihlerin stratejik olarak yanlış bildirilmesi \cite{ferragut2022scheduling} modelin dışında kalmaktadır.
\Sec{sec:invariance} içindeki iki EV'li örneğin gösterdiği gibi toplam olurluluk (A5) kapalı döngüde sağlanamayabilir.
Bu tür adımlarda ayrılış gereksinimi yalnızca Teorem~\ref{thm:invariance}(b) kapsamındaki EV'ler için garanti edilir; Lemma~\ref{S-lem:deficit} ise kalan eksik enerjiyi ortadan kaldırmadan sınırlar.
Gelecek boyunca ortak olurluluğu sertifikalayan bir katman, zamana bağlı kısıtlara ihtiyaç duyacak ve tek toplamlı izdüşüm yapısını kaybedecektir.

Trafoya ilişkin sonuç, A3 tek adım modelini ve A6'daki dışsal kapasite payını gerektirir; denetleyici A6'yı uygulamaz.
Lemma~\ref{S-lem:a3}, kullandığımız EV2Gym ayarları için A3'ü ispatlamaktadır; buna karşın gerçek bir batarya yönetim sistemi komut edilenden daha az güç çekebilir.
Eksik teslimat, bir EV'nin enerjisini sınırların varsaydığı değerin altına düşürür; bu nedenle donanım uygulamasında eşikler üzerinde bir pay ya da ölçülen enerjiden geri besleme gerekecektir; bu çözüm burada analiz edilmemiştir.
A1, EV2Gym'in üstel şarj gücü azalması için geçerlidir; ancak çok daha dik bir azalma bu varsayımı bozacaktır.
Model, gerilim limitleri ya da hat akışları olmaksızın tek bir trafonun arkasındaki tek bir istasyonu kapsar ve ek bağlantılı toplam kısıtları Önerme~\ref{prop:projection} kapsamının dışında kalacaktır.

Deneylerde EV2Gym'in kayıpsız muhasebesi kullanılmakta ve deşarj fiyatı şarj fiyatına eşit alınmaktadır.
Dolayısıyla V2G arbitrajı ne gidiş-dönüş kaybı ne de yıpranma maliyeti içerir; bu nedenle deşarjdan elde edilen kârlar iyimser olabilir ve duyarlılık çalışmaları bu etkiyi yalnızca tek bir istasyon boyutunda incelemektedir.
Tüm sonuçlar, EV2Gym istatistiklerinden üretilen sentetik varışlar ve Hollanda gün öncesi fiyatlarıyla yapılan simülasyonlardan elde edilmiştir; bazı geç akşam oturumları epizot ufkuna sığmamakta ve tahmin gürültüsü sentetiktir.
Şebeke büyüklükleri $N$ ile doğrusal olarak ölçeklendiğinden test istasyonlarının şarj ünitesi başına istatistikleri eğitimdekilerle örtüşmekte; bunları yalnızca stres düzeyi değiştirmektedir.

Bazı karşılaştırmalar tam olarak eşdeğer koşullarda yapılmamıştır.
EV-GNN, kendi algoritması ve ödülüyle, tek bir boyutta, her profil için ayrı olarak, çalıştırma başına tek bir yük ve PV çekilişiyle ve varsayılanından daha az adımla eğitilmektedir.
eMPC-HiGHS yeniden gerçeklememiz ufku içindeki kesin fiyatları ve varışları bilmekte, MLP yalnızca $N\le20$ için raporlanmakta ve GNN-DT çalıştırılmamıştır.
Optimumlar, tek adımlık bir geçiş bandı ve \%1'lik bir pay dışında kesin olan dış şarj gücü azalması satırlarını yuvarlama payıyla birlikte kullanmaktadır.
Ayrıca optimumlar çevrim içi bir denetleyicinin bilgi kısıtlarını göz ardı etmektedir; dolayısıyla her açığın bir kısmı geleceği bilmemenin bedelidir.

% ---------------------------------------------------------------------
\section{Sonuç}\label{sec:conclusion}
% ---------------------------------------------------------------------

Bu makale, bir CPO istasyonunda V2G şarj kontrolü için boyuttan bağımsız küme politikalarını olurluluğu koruyan bir güvenlik katmanıyla birleştirmiştir.
Katman, eğitim ve kullanım sırasında ham eylemleri şarj gücü azalmasını dikkate alan araç başına sınırlara ve iki yönlü trafo limitine $O(N\log N)$ sürede izdüşürmektedir.
Teorem~\ref{thm:invariance} iki katmanlı bir sonuç vermektedir.
Tek adım modeli ve şebeke kapasite payı koşulu altında trafo limiti her ham eylem dizisi için her adımda sağlanır ve kalış süresi boyunca olursuz adımla karşılaşmayan her EV kabul edilen ayrılış gereksinimini karşılar.

EV2Gym'de güvenlik katmanından geçirilerek çalıştırılan hiçbir epizotta trafo aşırı yüklemesi oluşmamıştır.
Ayrılışın bilindiği sınırı kullanan konformal olmayan her yöntemde, olursuz adım içermeyen hiçbir epizotta eksik enerjiyle ayrılan EV olmamıştır.
10 ve 20 şarj ünitesinde eğitilen dikkat tabanlı ve Deep Sets politikaları 200 şarj ünitesine aktarılabilmiştir.
Nokta kestirimlerinde dikkat tabanlı politika, 50 şarj ünitesinde eğitilmiş bir politikanın şarj ünitesi ve epizot başına 0,04\,\texteuro{} yakınında kalmıştır.
Bir mimari olarak dikkat, Deep Sets'e göre bir üstünlük göstermemiştir.
Aynı katmanla kullanılan bir fiyat eşiği kuralı, test ızgarasının her istasyon boyutunda dikkat tabanlı politikadan, 20 şarj ünitesinin altında ise Deep Sets politikasından daha düşük ortalama açığa sahip olmuştur.

Öğrenen politikalar olursuz adımlara yol açmış; istasyon boyutuna bağlı olarak dikkat tabanlı politikanın kabul edilen EV'lerinin \%7,4 ile \%34,1'i kalış süreleri boyunca bu tür bir adımla karşılaşmıştır.
Ayrılış sınırıyla dikkat tabanlı politika, sınır içermeyen bir varyanta ve ceza ile eğitilen ikizine göre daha fazla EV'yi gereksinimi karşılanmadan bırakmıştır.
Cezalı kâr açısından bu ikizden daha yüksek puan almış; 200 şarj ünitesindeki fark ise sonuçsuz kalmıştır.
\TODO{Yazarlar Tartışma bölümünü yeniden yazdıktan sonra bu sonuçları doğrular ya da yeniden yazar}.

Bir CPO açısından üç sınırlılık en çok önem taşımaktadır.
Birincisi, ayrılış sonucu bildirilmiş gereksinimler ve ayrılış zamanları ya da yalnızca marjinal kapsamaya sahip kalibre edilmiş bir ayrılış sınırı gerektirir.
İkincisi, trafoya ilişkin sonuç her EV'nin komut edilen gücü çektiğini varsayar; gerçek batarya yönetim sistemleri ise daha az güç çekebilir.
Üçüncüsü, kanıtlar öğrenen yöntem başına iki eğitim tohumuyla ve kayıpsız ekonomik modelle yapılan simülasyonlardan gelmektedir.
Gelecekteki çalışmalar olursuz adımları azaltan kapasiteyi dikkate alan bir alt sınır, donanım için ölçülen enerjiden geri besleme ve doğrusallaştırılmış gerilim limitleriyle birden fazla trafo ekleyebilir.
Kaydedilmiş şarj oturumlarıyla yapılacak testler, model ile uygulama arasında kalan farkı daraltabilecektir.



% ---------------------------------------------------------------------
% Makale Bilgi Formu (SAUCIS son kısmı; Kaynaklar'dan önce gelir)
% Maddeler SAUCIS LaTeX şablonunu izler. Anonim anahtarı adları ve
% bağlantıları gizler, ancak beyanları korur.
% ---------------------------------------------------------------------
\section*{Makale Bilgi Formu}

\textbf{Yazarların Katkıları}\\
\ifanonymous
Çift kör değerlendirme için kaldırılmıştır.
\else
\TODO{Her yazarın katkısını belirtin; örneğin ``Yazar 1: kavramsallaştırma, yöntem, yazılım, yazım; Yazar 2: danışmanlık, gözden geçirme ve düzenleme''; yenilenen DergiPark sistemi ayrıca bir katkı rolü beyanı da isteyebilir}.
\fi

\textbf{Teşekkür}\\
\ifanonymous
Çift kör değerlendirme için kaldırılmıştır.
\else
\TODO{Yazarlık ölçütlerini karşılamayan katkı sahiplerine teşekkür edin ya da ``Uygulanamaz'' yazın}.
\fi
\TODO{Teşekkür (ya da Yöntem) bölümü için yapay zekâ kullanımına ilişkin cümle, yazarlar tarafından yazılacaktır; SAUCIS bu beyanın Yapay Zekâ Beyanı'nın yanı sıra Yöntem ya da Teşekkür bölümünde de yapılmasını şart koşmaktadır}.

\textbf{Çıkar Çatışması}\\
Yazarlar, bu makalenin yayımlanmasıyla ilgili herhangi bir çıkar çatışması bulunmadığını beyan eder.
\TODO{Bu beyanı onaylayın}.

\textbf{Destek/Destekleyen Kuruluşlar}\\
\TODO{Her mali destek kaynağını proje numaralarıyla belirtin ya da herhangi bir mali destek alınmadığını ifade edin}.

\textbf{Etik Onay}\\
\TODO{Aşağıdaki ifadeyi onaylayın}.
Bu çalışma; insan katılımcı, yazarlarca toplanmış insan verisi ya da hayvan içermeyen bir simülasyon çalışması olduğundan etik kurul onayı gerekmemiştir.
Yazarlar, araştırma ve yayın etiğine uyduklarını beyan eder.

\textbf{Erişilebilirlik}\\
Kod, senaryo yapılandırmaları, test tekrarları ve analiz betikleri herkese açık bir depoda yayımlanacaktır.
\ifanonymous
Bağlantı, çift kör değerlendirme nedeniyle paylaşılmamıştır.
\else
\TODO{Depo bağlantısı, lisans ve arşivlenmiş sürüm}.
\fi
Deneylerde \texttt{55ef368c} commit'indeki açık EV2Gym kodu, kaynak kodunda değişiklik yapılmadan kullanılmaktadır.

\textbf{Yapay Zekâ Beyanı}\\
\TODO{Yazarlar tarafından yazılacaktır. Yapay zekâ katkılarının olgusal bir kaydı \texttt{review/\allowbreak ai\_contribution\_record.md} dosyasında tutulmaktadır}.

\textbf{İntihal Beyanı}\\
Bu makale iThenticate\texttrademark{} ile taranmıştır.


% ---------------------------------------------------------------------
% Terim Sözlüğü (yalnızca Türkçe inceleme kopyasında; numarasız ek)
% ---------------------------------------------------------------------
\clearpage
\section*{Terim Sözlüğü}
\addcontentsline{toc}{section}{Terim Sözlüğü}

Bu çeviride kullanılan terimler aşağıda listelenmiştir (İngilizce $\rightarrow$ Türkçe).
Matematikte, tablolarda ve şekillerde kullanılan kısaltmalar (RL, EV, CPO, V2G, LLF, PPO vb.) İngilizce biçimleriyle korunmuştur.
Karşılaştırılan yöntemlerin etiketleri (Attention+penalty, Attention+PPO-Lagrangian, Deep Sets+safety layer, eMPC-HiGHS vb.) şekil açıklamalarıyla tutarlı olması için İngilizce bırakılmıştır.

{\small
\setlength{\LTpre}{2pt}\setlength{\LTpost}{0pt}
\begin{longtable}{@{}>{\raggedright}p{0.44\linewidth}>{\raggedright\arraybackslash}p{0.50\linewidth}@{}}
\toprule
İngilizce & Türkçe \\
\midrule
\endfirsthead
\toprule
İngilizce & Türkçe \\
\midrule
\endhead
\bottomrule
\endfoot
actor / critic & aktör / kritik \\
admission control & kabul kontrolü \\
admitted requirement & kabul edilen gereksinim \\
attention policy & dikkat tabanlı politika \\
baseline & karşılaştırma yöntemi (baseline) \\
charge point operator (CPO) & şarj noktası işletmecisi (CPO) \\
charger & şarj ünitesi \\
charging taper & şarj gücü azalması (taper) \\
common shift & ortak kaydırma \\
conformal prediction; split-conformal & konformal tahmin; bölünmüş konformal \\
cost of the guarantee & garantinin maliyeti \\
coverage & kapsama \\
deficit bound & açık sınırı \\
departure (energy) requirement & ayrılış (enerji) gereksinimi \\
deployment & kullanım (devreye alma) aşaması \\
electric vehicle (EV) & elektrikli araç (EV) \\
EV left unserved / short & gereksinimi karşılanmadan bırakılan EV \\
endpoint (primary) & (birincil) sonlanım noktası \\
energy not served & karşılanmayan enerji \\
episode & epizot (episode) \\
fallback rule & yedek kural \\
feasibility; feasible; infeasible step & olurluluk; olurlu; olursuz adım \\
grid headroom & şebeke kapasite payı \\
hard / soft (constraint, optimum) & katı / yumuşak (kısıt, optimum) \\
inconclusive (comparison) & sonuçsuz (karşılaştırma) \\
inflexible load & esnek olmayan yük \\
Lagrangian & Lagrange \\
latest-start trajectory & en geç başlangıç yörüngesi \\
laxity & esneklik payı (laxity) \\
least-laxity-first (LLF) & en az esneklik payı önce (LLF) \\
lower / upper bound & alt / üst sınır \\
mask; key-padding mask & maske; anahtar dolgu maskesi \\
model predictive control (MPC) & model öngörülü kontrol (MPC) \\
optimality gap & optimallik açığı \\
overload (energy) & aşırı yük (enerjisi) \\
penalised profit & cezalı kâr \\
penalty & ceza \\
perfect-information optimum & tam bilgili optimum \\
policy & politika \\
projection & izdüşüm (projeksiyon) \\
raw action; commanded power; realised power & ham eylem; komut edilen güç; gerçekleşen güç \\
reachable energy & erişilebilir enerji \\
reinforcement learning (RL) & pekiştirmeli öğrenme (RL) \\
rank-biserial correlation & sıra-çift serili korelasyon \\
replay (test replay) & tekrar (test tekrarı, replay) \\
request pipeline & talep işleme hattı \\
reward & ödül \\
rollout & veri toplama turu (rollout) \\
safety layer & güvenlik katmanı \\
seed & tohum (seed) \\
seed-averaged policy & tohum ortalamalı politika \\
set policy & küme politikası \\
shortfall & eksik enerji (eksiklik) \\
size-agnostic & boyuttan bağımsız \\
smoothed LLF (sLLF) & yumuşatılmış LLF (sLLF) \\
step & adım \\
stress level & stres düzeyi \\
token & belirteç (token) \\
transformer (power) & transformatör (trafo) \\
transformer encoder (neural) & Transformer kodlayıcı \\
transformer limit; transformer rating & trafo limiti; trafo anma gücü \\
vehicle-to-grid (V2G) & araçtan şebekeye (V2G) \\
zero-shot transfer & sıfır-atış aktarım (yeniden eğitim olmadan) \\
\end{longtable}
}


\IfFileExists{IEEEtran.bst}{\bibliographystyle{IEEEtran}}{\bibliographystyle{ieeetr}}
\bibliography{refs}

\end{document}
